\documentclass[10pt]{article}

\usepackage{tmlr}
\usepackage{amsmath,amssymb}
\usepackage{booktabs}
\usepackage{graphicx}
\usepackage{xcolor}
\usepackage{url}
\usepackage{placeins}
\graphicspath{{figs/}}

\newcommand{\delp}{DeLP}
\newcommand{\acc}{\mathrm{acc}}
\newcommand{\own}{p_{\mathrm{own}}}

\title{What an Active-Dendrite Continual Learner Actually Learns: A Dead-Neuron Artifact, a Granularity Result That Did Not Survive Its Own Controls, and a Repair That Is a Damping Factor}
\author{\name Anonymous Author \email anonymous@example.com \\ \addr Anonymous Institution}

\begin{document}

\maketitle

\begin{abstract}
Active-dendrite architectures route each input through a competition among
dendritic segments, one of which is trained per task; usage-localized elastic
weight consolidation (\delp{}) protects what has been learned by scaling
per-neuron Fisher anchors with segment usage. Both mechanisms rest on an
implicit assumption that has, to our knowledge, never been measured: that the
segment--task mapping constitutes a \emph{specialization} that protection can
leverage at sub-neuron granularity. We audit this assumption on
Permuted-MNIST and Permuted-Fashion and report eight findings, five of them
negative, that together invert the interpretation of a repair we had
originally reported as causal. \textbf{First}, a granularity ladder of anchors
(plain, neuron-resolved, segment-resolved with attribution-tracked reference
states) shows that the benefit of usage localization measured on the learning term is
realized at neuron level in every cell we ran: segment-resolved anchors differ from
neuron-resolved anchors by no more than the resolution of a 10-seed paired test
across every combination of two datasets, two task
counts ($T{=}10,20$), and two values of the elastic weight
$\lambda$ ($|\Delta\ell|\le 0.24$, all
$n{=}10$ paired tests n.s.), and remain so under FIFO anchor budgets
($|\Delta\ell|\le 0.44$ across eight budget cells). \textbf{Second}, the audit
explains why: no task-dominant gating is \emph{detectable} by the standard
probability that a neuron's dominant segment is its own task's segment starts
each phase at chance and is pushed \emph{below} chance by that phase's own
training ($T{=}20$: $0.05 \to 0.01$), and Fisher mass is diffuse across
segments (top-1 share $11$--$18\%$). Because that instrument is itself unreliable
in this family (below), we read this as an absence of measurement rather than a
measured absence. The audit also uncovers a measurement
artifact relevant to the active-dendrite literature: under production sampling
budgets $56$--$63\%$ of neurons are never activated, and naive argmax
instrumentation over all neurons fabricates a strong ``segment-0 dominance''
($\approx 60\%$) that vanishes under active-only filtering.
\textbf{Third}, a constant bias $b$ on the own-segment dendritic input appears
to repair this: it raises the learning term in every cell of a
$\{$dataset$\}\times\{$depth$\}\times\{$granularity$\}$ grid
($\Delta\ell>0$ in $22/22$ dose cells, significant in $16/22$ at $n{=}10$
paired; $\Delta\ell$ up to $+4.6$, $t(9){=}8.1$). \textbf{Fourth}, however, the
repair does not transfer. On a convolutional active-dendrite backbone the
dose--response collapses to noise ($0/12$ cells survive correction) and final accuracy
reverses (down to $-2.29$, $t{=}-3.1$); a same-architecture control separates
the candidate causes by showing that the predicted link between mortality and repair gain does not hold
(across eight configurations of a width$\times$normalization factorial scan the two are
uncorrelated, $r=-0.11$; GroupNorm raises mortality to $76.5\%$/$85.5\%$ while collapsing the gain to
$17\%$/$26\%$ of its bare-MLP value), and a $2\times 2$ width $\times$
normalization scan ($320$ runs) confirms the attribution --- normalization
collapses the gain at both widths, while widening the trunk $7.8\times$
\emph{increases} it. \textbf{Fifth}, a mechanism probe re-qualifies what the
repair does: the winning segment's gate value is a constant of the
initialization ($0.565$, already $0.564$ before any training), identical under
every alternative context; and we trace why the context pathway is unused to a
defect in the protection protocol itself, whose per-task gradient mask freezes
$90\%$ of the gate matrix at its random initialization while the forward pass
reads whole columns --- so the context pathway is not merely unused but
untrainable; at $b{=}10$ the own segment always wins with
$\sigma(10+u)\approx 1$, making the network context-degenerate, and a
four-arm control (base / constant-gate / repaired / gate-deleted) shows that
gate \emph{deletion} recovers most of the gain on MNIST ($76\%$, against an
anchor-weighting sensitivity of $\pm3.9$ points) and is numerically above the
repair on Fashion-MNIST, a difference that is not significant at $n{=}10$. The honest summary is that we measured the size, the cost, and
the boundary of the repair precisely, and located its mechanism only far enough
to know that it is not the one we first assumed: the bias does not restore
routing, it removes a constant damping factor that a never-used gate had left
in place. Sixth, the mechanism result has an uncomfortable implication for our
own second finding, and we state it rather than leave it for a reviewer: if
the gate's value does not depend on which segment wins, then ``dominant
segment'' is not a description of a routing decision but of an argmax over
near-degenerate rows. We check what survives --- the granularity ladder
compares anchor constructions and never uses a routing statistic, the
dead-neuron artifact is a property of argmax on a zero block, and the
mechanism and ablation measure outputs and interventions --- and adopt the
weaker formulation that check licenses: we cannot claim that specialization
is actively violated as a property of learned computation, only that no
evidence of it is detectable by an instrument that is itself unreliable in
this family. Segment--task specialization in this architecture family is
therefore not a weak tendency but an assumption with no support that survives
an instrument we can trust --- and what the repair buys is not compensation for
that absence but the removal of a constant damping factor. Seventh, the
controls we added after the first review of this line of work change three of
our claims. A positive control shows the dominant-segment instrument is
sensitive: when specialization is forced by construction, the instrument reads
$1.000$ against $0.055$ for the base network, so the negative readings are real.
A specificity control shows the bias acts on whichever segment the network is
using, not specifically on the task's own, and that a pure constant rescaling of
the trunk does better still. And a construction-diversity control shows the
granularity null was an artefact of our self-limiting reference rule: a
refreshed-reference segment-resolved anchor gains about two and a half points of learning
term on MNIST while losing about twenty points of final accuracy, so we
withdraw the general claim that sub-neuron granularity is worthless. Eighth, the
method we audit is not competitive with contemporary alternatives: a replay buffer
with logit distillation reaches $86.4$ against DeLP's $64.0$ on Permuted-MNIST
and $75.3$ against $42.5$ on Permuted-Fashion, with markedly less forgetting in
both, in all twelve paired comparisons significant at $p<10^{-6}$. This paper is
therefore an argument that the dendritic-segment account of continual learning
needs to be rethought from the trunk upward, not a defence of the method it audits.
\end{abstract}

\noindent\textbf{Keywords:} continual learning , elastic weight consolidation , active dendrites
, dendritic segments , catastrophic forgetting , interpretability
audit , reproducibility , negative results , experience replay
, dark experience replay , equivalence testing

\section{Introduction}
\label{sec:intro}

Biological neurons integrate thousands of synaptic inputs across morphologically
distinct dendritic compartments, and single-neuron studies have long argued
that these compartments act as quasi-independent computational subunits
\citep{poirazi2003,london2005,poirazi2020}. An influential line of work has
translated this principle into artificial networks: context signals select
one of several dendritic \emph{segments} per neuron, a winner-take-all
competition among segments determines the nonlinear contribution to the
soma, and \emph{only the selected segment receives gradient} on a given task
\citep{iyer2022}. The architectural promise is a division of labor: if task
$k$'s computation flows predominantly through segment $k$, then tasks
naturally occupy disjoint sub-network state, and interference between tasks
has a structural bottleneck.

Continual-learning methods built on dendritic architectures inherit this
promise, often silently. \delp{} (usage-localized elastic weight
consolidation) protects per-neuron parameter blocks with Fisher anchors
scaled by how heavily each neuron was used in past tasks
\citep{zhou2026delp,kirkpatrick2017}. Its protection unit is the neuron; the
segment structure below the neuron is used only to define usage. A natural
proposal --- one we initially expected to succeed --- is to push protection
\emph{down} to the segment level: if tasks really live in different segments,
anchors that track each segment's own reference state should protect more
precisely and drag less on future learning, echoing the granularity
intuition that finer protection localizes interference better.

This paper reports what actually happens when that assumption is tested
rather than assumed. The narrative is a sequence of three failures that
invert the interpretation of the last one, and we report it in that order
because the ordering is the finding: an apparent success in the fourth
contribution below is not a success, and we say so.

Our contributions are sixfold, and each arose from the failure of the
previous one:

\begin{enumerate}
\item \textbf{A granularity audit with a null result at segment level
      (Sec.~\ref{sec:granularity}).} We implement segment-resolved anchors
      with \emph{attribution-tracked reference states} --- each anchor tracks
      where its own segment's parameters drift --- and compare a three-rung
      ladder (plain, neuron-resolved, segment-resolved) across two datasets,
      two sequence lengths, and two values of $\lambda$ with $n{=}10$ paired seeds. The
      entire benefit of usage localization is realized in the
      plain$\to$neuron step ($\Delta\ell$ between $-2.8$ and $-5.4$,
      $|t(9)|$ up to $22$); the neuron$\to$segment step adds nothing
      ($|\Delta\ell|\le 0.24$, all n.s.), under unbounded stacks and under
      FIFO anchor budgets alike. The null is consistent with two properties of the
      construction we chose, rather than merely observed: attribution is
      \emph{self-limiting} (the anchor's own
      penalty pins the weights, so the tracked reference never moves) and
      segment-level parameter ownership lacks a separable carrier.
      \textbf{This null does not survive a construction that is not}
      \textbf{self-limiting} (Sec.~\ref{sec:ctrl-anchor}): with refreshed
      per-segment references the same segment-resolved anchor gains
      about two and a half points of learning term on MNIST, at a cost of roughly
      twenty points of final accuracy. We therefore withdraw the general
      claim that sub-neuron protection granularity is worthless: what the
      control shows is that the direction of the effect is set by the
      reference-update rule, not by the granularity itself.
\item \textbf{A measurement artifact with literature-wide relevance
      (Sec.~\ref{sec:audit}).} Auditing \emph{why} the null occurs, we find
      that under production sampling budgets $56$--$63\%$ of neurons are
      never activated, and that naive dominant-segment instrumentation
      (argmax over segment activations of \emph{all} neurons) attributes
      every dead neuron to segment 0 --- fabricating an apparent
      ``segment-0 dominance'' of $\approx 60\%$ that vanishes under
      active-only filtering. Among active neurons the picture is
      \emph{anti-specialization} \emph{as measured}: each phase's own training pushes the
      own-segment win rate from chance to below chance
      ($T{=}20$: $0.05 \to 0.01$), and anchor Fisher mass is diffuse across
      segments (top-1 share $11$--$18\%$). Any study that measures segment
      dominance by argmax over segment activations \emph{and} does not filter
      inactive units is at risk of the same artifact; we did not audit the
      instrumentation used in other work on this architecture.
\item \textbf{An apparent causal repair (Sec.~\ref{sec:causal}).} If missing
      specialization is load-bearing rather than cosmetic, restoring it
      should help. It does: adding a constant bias $b$ to the own-segment
      dendritic input during training, evaluation, and anchor
      construction raises the learning term in every cell of a
      $\{$dataset$\}\times\{$depth$\}\times\{$granularity$\}$ grid ---
      $\Delta\ell$ positive in $22/22$ dose cells, significant in $16/22$
      ($n{=}10$ paired) --- with final accuracy significant simultaneously
      in $16/22$ and in every grid cell at some dose ($\Delta\ell$ up to
      $+4.6$ at $T{=}20$, $t(9){=}8.1$); forgetting is \emph{reduced} at
      $T{=}10$ ($\Delta\mathrm{BWT}$ up to $+1.7$, $t(9)$ up to $10.7$) and
      mildly traded back at $T{=}20$ (significant in $3/10$ cells)
      (Sec.~\ref{sec:causal}). The response is saturated and nontrivially
      mediated: the own-segment win rate saturates at $13$--$19\%$
      regardless of dose. Signed-drift measurement shows this is not
      adversarial compensation --- the network never drifts against the bias
      --- but a scale effect: the raw own gate spans $|u| \le 0.32$, at least
      three times below any tested dose, so the bias forces the same
      training-time routing at every dose, so only a minority of neuron-tasks
      flip to the own-segment pathway --- and that partial restoration
      already yields most of the gain.
\item \textbf{The repair does not transfer, and the reason is not gate death
      (Sec.~\ref{sec:conv}).} The third contribution is the one we most
      wanted to be wrong about. On a convolutional active-dendrite backbone
      ($160$ runs) the dose--response collapses to noise ($2/12$ cells
significant uncorrected, $0/12$ after Holm correction) and final accuracy \emph{reverses} (Fashion $T{=}10$:
      $\Delta\mathrm{avg}=-2.29$, $t{=}-3.1$). The natural explanation ---
      the repair revived dying gates, and conv gates do not die --- is
      testable, and we test it: inserting GroupNorm into the MLP raises
      mortality to $76.5\%$/$85.5\%$ while collapsing the repair gain to
      $17\%$/$26\%$ of its bare-MLP value. Mortality and gain
      do not track each other (the highest-mortality configuration has one of the
      smallest gains;
      zero mortality with no gain), so gate death is neither necessary nor
      sufficient. A $2\times 2$ width $\times$ normalization scan
      ($320$ runs) removes the one confound we had to disclose: normalization
      collapses the gain at both widths, while widening the trunk $7.8\times$
      \emph{increases} it.
\item \textbf{What the repair actually repairs (Sec.~\ref{sec:mech}).} A
      mechanism probe re-qualifies the third contribution. The winning
      segment's gate value is a constant of the initialization ($0.565$,
      already $0.564$ before training), identical to four decimals under
      every alternative context, and the dendritic columns stay mutually
      orthogonal at their random initialization (mean pairwise cosine
      $0.003$, cosine to the initialization $0.9995$): the context pathway
      is present, protected per task, and \emph{unused as a modulator}.
      We then trace the cause to the protection protocol rather than to
      the argmax: its per-task gradient mask leaves every off-diagonal
      block of the gate matrix at exactly its random initialization
      ($0.0$ relative displacement in all 90 blocks, every seed; $n=5$
      seeds per dataset, independently re-run on a second machine), while
      the forward pass reads whole columns, so $90\%$ of the context
      pathway never receives a gradient at all. At
      $b{=}10$ the own segment always wins with $\sigma(10+u)\approx 1$,
      which makes the network context-degenerate (swapping the context
      one-hot changes $0.007$--$0.021\%$ of predictions, against
      $14.8$--$30.1\%$ for the base network). A four-arm control (base /
      constant-gate / repaired / gate-deleted, $n{=}10$ paired) settles
      what the bias buys: pinning the gate at the base system's measured
      constant reproduces the base system ($63.62$ vs.\ $63.86$ on MNIST,
      n.s.), while \emph{deleting} the gate reproduces most of the repair's
      gain on MNIST ($67.50$ vs.\ $68.67$) and is numerically above it on
      Fashion-MNIST ($44.27$ vs.\ $42.89$, a difference that is not significant
      at $n{=}10$). The bias does not restore
      routing; it removes a constant damping factor that a never-used gate
      had left in place.
\item \textbf{A self-interrogation we should have run first
      (Sec.~\ref{sec:pseudo}).} The mechanism result has an uncomfortable
      implication for our own second contribution, and we state it rather
      than leave it for a reviewer: if the gate's value does not depend on
      which segment wins, then ``dominant segment'' is not a description
      of a routing decision but of an argmax over near-degenerate rows.
      We check what survives. The granularity ladder
      (Sec.~\ref{sec:granularity}) compares anchor constructions and never
      uses a routing statistic; the dead-neuron artifact is a property of
      argmax on a zero block; the mechanism and ablation measure outputs
      and interventions. All three stand. What does \emph{not} survive is
      the strong form of our second contribution: we cannot claim that
      specialization is \emph{actively violated} as a property of learned
      computation, only that no evidence of it is detectable by an
      instrument that is itself unreliable here. We adopt the weaker
      formulation throughout.
\end{enumerate}

\paragraph{Scope relative to a companion study.} A separate
paper in this line asks a different question about the same system: how
large the protection set should be. It treats the anchor stack as the
object of study, bounds the accumulated Fisher mass, and derives a
log-linear law relating the learning term to the protected-mass horizon
($\ell \approx 84 - 11\ln h$; per-seed $R^2 \ge 0.98$; replicated at a
second task count and on Split-CIFAR-100 with a ResNet-18 backbone).
That question and this one are compatible but independent. Bounding the
stack answers how much protection to keep, not whether finer protection
would be worth having; and this paper's answer --- that protection
granularity below the neuron buys nothing, because the sub-neuron
structure it would exploit is not there --- implies that a bounded stack
of neuron-resolved anchors is the right object to bound. Readers who
want the account of \emph{how much} to protect should consult that study;
this paper is about \emph{what} there is to protect.

Together the six results reframe segment specialization in this
architecture family: not a weak tendency to be exploited, but an
assumption with no support that survives a corrected instrument --- and
whose absence a one-parameter intervention compensates for rather than
cures. The last of these results is the one we would most like to have
run first, and the reason we did not is worth stating: the apparent
repair was strong, consistent, and saturating, and that is exactly what
made it seem unnecessary to ask what the gate was doing. All experiments are on
Permuted-MNIST and
Permuted-Fashion, on one-hidden-layer MLPs except for the convolutional and
width-variant arms of Secs.~\ref{sec:conv}--\ref{sec:gnmlp}, at $T\in\{10,20\}$
with paired-seed statistics; code,
per-run records, analysis scripts, and the pre-registrations are public
(Data Availability).

% =====================================================================
\section{Background and setup}
\label{sec:setup}

\subsection{Active-dendrite network}
\label{sec:adnet}

\label{sec:adnet}

We use the feed-forward active-dendrite network of Iyer et al.\
\citep{iyer2022} in the protected variant established in our prior work
\citep{zhou2026delp}. A one-hidden-layer network of $N_h$ neurons receives a
one-hot task context $c_t \in \{0,1\}^{T}$; each neuron $j$ carries $T$
dendritic segments, and segment $k$ of neuron $j$ computes
$u_{j,k} = \sigma(w_{(j,k)}^\top x + b_{j,k} + v_{k}^\top c_t)$, where
$v \in \mathbb{R}^{T \times (n_{seg} N_h)}$ are the context-to-dendrite
weights --- \emph{column} $k$ of $v$ gates segment $k$ for every neuron.
The soma nonlinearity applies a $k$-winners-take-all (kWTA) selection over
$h_j = f_j(x)\, \sigma(\max_k u_{j,k})$: the $k$ largest somal pre-activities
pass; the rest are set to zero. Crucially, task $t$'s gradient updates are
masked to the parameters of segment $t$ (``protect\_grads''): the
architecture's design intent is that task $t$ trains segment $t$.

\subsection{\delp{}: usage-localized EWC and the identity decomposition}
\label{sec:delp}

\label{sec:delp}

\delp{} \citep{zhou2026delp} protects the feed-forward and output blocks with
per-task anchors, each an elastic penalty
$\frac{\lambda}{2}\sum_p F_p\, (W_p - \theta_p)^2$ whose Fisher term $F$ is
\emph{usage-localized}: per-neuron usage $o_j \in [0,1]$, the frequency with
which neuron $j$ is active under task $t$'s data, scales the anchor rows of
the hidden weights and columns of the output weights, so that rarely-used
neurons are protected lightly and remain available to future tasks. Writing
$\bar a_T$ for final average accuracy, the decomposition
\begin{equation}
\bar a_T \;=\; \ell \;+\; \tfrac{T-1}{T}\,\mathrm{BWT},
\label{eq:identity}
\end{equation}
separates what each task achieves when it is learned ($\ell$, the learning
term: the mean over tasks of the accuracy on a task immediately after that
task has been trained) from what earlier tasks retain (BWT $\le 0$); we
report all three, with paired $t$-statistics over $n{=}10$ seeds sharing
permutation streams. Where a table quotes an accuracy on the \emph{final}
task alone, it says so explicitly (Table~\ref{tab:refs}).

\subsection{Protocol}
\label{sec:protocol}

\label{sec:protocol}

Benchmarks are Permuted-MNIST and Permuted-Fashion \citep{parisi2019,xiao2017} with
$T\in\{10,20\}$ tasks, one hidden layer of $256$ neurons, $2$ epochs per
task, Adam (lr $10^{-3}$), Fisher/usage samples of $300$/$500$ examples, and
$\lambda \in \{1000, 5000\}$ calibrated in prior work. Seeds are paired
across arms (identical permutation streams and init); all comparisons are
paired $t$-tests with $n{=}10$, and we mark $|t|>2$ as significant. All
hypotheses in Secs.~\ref{sec:granularity} and \ref{sec:causal} were
preregistered in writing before the corresponding data were collected.

% =====================================================================



% ===== the bare-MLP arm 主体 =====


\section{Does segment-resolved protection help? A granularity ladder}
\label{sec:granularity}

\subsection{Segment-resolved anchors with attribution-tracked references}
\label{sec:branchdesign}

A segment-resolved anchor must differ from a neuron-resolved one in more
than bookkeeping: since $\sum_k o_{t,j,k} = o_{t,j}$ (segment usages
partition neuron usage), splitting a neuron's single Fisher scalar across
segments while sharing one reference state $\theta$ leaves the penalty
algebraically identical to neuron level. We therefore give each segment
attribution of parameter drift: anchor $a$ (task $t$, neuron $j$) maintains
per-segment reference states $\theta^{(a,k)}$, and after each subsequent
task $t'$ the realized parameter displacement is partitioned across segments
by the usage-driven attribution weights of $t'$'s own anchors; the penalty
charges anchor $a$ only for the drift not explained by interlopers in its
own segment. The design intends that an anchor no longer litigates, on
shared weights, drift caused by tasks that never touched its segment.
Segment-level Fisher mass $m_{t,j,k}$ and dominant-segment histograms are
instrumented per anchor.

\subsection{Result: the ladder is flat above neuron level}

Table~\ref{tab:gran} shows the ladder. The plain$\to$neuron step is large
and universal: at $\lambda{=}1000$, $T{=}10$ MNIST the learning term gains
$+4.66$ ($t(9){=}22.3$); across all six cells of
$\{$dataset$\}\times\{$depth$\}\times\{\lambda\}$, $\Delta\ell$ ranges from
$-2.78$ to $-5.44$ in favor of neuron-level localization with
$|t(9)| \in [10.4, 22.3]$. The neuron$\to$segment step is empty on the learning axis: the
largest deviation across all twelve pairings is $|\Delta\ell| = 0.24$ on the learning
term and $|\Delta\mathrm{BWT}| = 0.43$ on forgetting,
never significant (largest $|t| = 1.3$); BWT differences are within
$0.43$ points. Two of those BWT differences are nominally significant
(effect sizes below $0.5$ points, $|t| \le 4.0$) and we report them rather
than suppress them: they are marginal forgetting effects, not evidence for
a granularity effect, since they sit on the stability side rather than the
learning side that the pre-registered hypothesis was about.

\begin{table}[t]
\centering\footnotesize\setlength{\tabcolsep}{4.5pt}
\caption{Granularity ladder (unbounded anchor stacks, $n{=}10$ paired
seeds). $\ell$/avg/BWT are absolute means; the right block is the paired
difference vs.\ the neuron-resolved arm. All plain$-$neuron differences are
significant at $|t|>10$. For branch vs.\ neuron the learning term is
indistinguishable in all six cells (largest $|t| = 1.3$ on $\Delta\ell$);
the only significant branch$-$neuron differences are marginal forgetting
effects (largest $|t| = 4.0$, $\Delta\mathrm{BWT} \le 0.43$ points), which
we do not read as a granularity effect.}
\label{tab:gran}
\begin{tabular}{llrrrr rr}
\toprule
 & & \multicolumn{3}{c}{absolute means} & \multicolumn{3}{c}{paired vs.\ neuron} \\
\cmidrule(lr){3-5}\cmidrule(lr){6-8}
cell & arm & $\ell$ & avg & BWT & $\Delta\ell$ & $\Delta$avg & $\Delta$BWT \\
\midrule
\multicolumn{8}{l}{\emph{Permuted-MNIST, $\lambda{=}1000$}}\\
$T{=}10$ & plain  & 69.63 & 61.25 & $-9.31$ & $-4.66^{*}$ & $-2.61^{*}$ & $+2.26^{*}$ \\
$T{=}10$ & neuron & 74.28 & 63.86 & $-11.58$ & --- & --- & --- \\
$T{=}10$ & branch & 74.21 & 64.13 & $-11.20$ & $-0.08$ & $+0.27$ & $+0.38^{*}$ \\
$T{=}20$ & plain  & 46.35 & 37.30 & $-9.53$ & $-5.44^{*}$ & $-1.35^{*}$ & $+4.31^{*}$ \\
$T{=}20$ & neuron & 51.79 & 38.64 & $-13.84$ & --- & --- & --- \\
$T{=}20$ & branch & 51.98 & 39.03 & $-13.63$ & $+0.19$ & $+0.39^{*}$ & $+0.21$ \\
\midrule
\multicolumn{8}{l}{\emph{Permuted-MNIST, $\lambda{=}5000$}}\\
$T{=}10$ & plain  & 54.41 & 51.69 & $-3.03$ & $-4.78^{*}$ & $-2.80^{*}$ & $+2.19^{*}$ \\
$T{=}10$ & neuron & 59.19 & 54.49 & $-5.22$ & --- & --- & --- \\
$T{=}10$ & branch & 59.21 & 54.71 & $-5.01$ & $+0.03$ & $+0.22$ & $+0.21^{*}$ \\
$T{=}20$ & plain  & 33.97 & 31.68 & $-2.40$ & $-3.72^{*}$ & $-1.74^{*}$ & $+2.08^{*}$ \\
$T{=}20$ & neuron & 37.68 & 33.42 & $-4.49$ & --- & --- & --- \\
$T{=}20$ & branch & 37.59 & 33.39 & $-4.42$ & $-0.09$ & $-0.03$ & $+0.07$ \\
\midrule
\multicolumn{8}{l}{\emph{Permuted-Fashion, $\lambda{=}5000$}}\\
$T{=}10$ & plain  & 38.53 & 33.47 & $-5.62$ & $-3.76^{*}$ & $-1.65^{*}$ & $+2.35^{*}$ \\
$T{=}10$ & neuron & 42.29 & 35.12 & $-7.96$ & --- & --- & --- \\
$T{=}10$ & branch & 42.05 & 35.27 & $-7.53$ & $-0.24$ & $+0.15$ & $+0.43$ \\
$T{=}20$ & plain  & 26.56 & 22.75 & $-4.01$ & $-2.78^{*}$ & $-1.22^{*}$ & $+1.65^{*}$ \\
$T{=}20$ & neuron & 29.35 & 23.97 & $-5.66$ & --- & --- & --- \\
$T{=}20$ & branch & 29.35 & 24.14 & $-5.49$ & $+0.01$ & $+0.17$ & $+0.17$ \\
\bottomrule
\end{tabular}
\end{table}

Single-anchor references (Table~\ref{tab:refs}) frame the arms; their
$\lambda$ values are the calibration carried over from our companion
pipeline, not re-tuned per benchmark. Plain EWC buys $\ell$ at
$3$--$5\times$ the forgetting; SI ($\lambda{=}100$) is over-stable (its
$\ell_{\mathrm{final}}{=}13.2$ at $T{=}10$ shows it barely learns the later
tasks at all, even though its mean learning term $\ell{=}40.0$ is not
unremarkable); MAS ($\lambda{=}1$) maximizes the learning term
($\ell \ge 85.7$, the highest of all arms) yet posts catastrophic BWT
($-32.6$ at $T{=}20$) --- the regime-dependence of warrant documented
in our companion work extends from Split-CIFAR-100 to the MLP setting. One
comparison the table makes possible is worth stating explicitly: at
$T{=}10$, the repaired \mbox{DeLP} arm ($b{=}10$; avg $68.7$, BWT $-9.9$)
dominates MAS (avg $66.6$, BWT $-21.2$) on both summary terms --- MAS
buys its higher learning term with exactly the forgetting that the repair
removes.

\begin{table}[t]
\centering\small
\caption{Single-anchor references (Permuted-MNIST, $n{=}10$). $\ell$ is the
learning term of Eq.~\eqref{eq:identity} (mean over tasks of the accuracy
right after each task is learned); $\ell_{\mathrm{final}}$ is the accuracy on
the last task alone and is reported separately because for these two arms it
differs substantially from $\ell$ --- SI in particular barely learns the later
tasks at all, which is visible only in $\ell_{\mathrm{final}}$.}
\label{tab:refs}
\begin{tabular}{llrrrr}
\toprule
 & method ($\lambda$) & $\ell$ & $\ell_{\mathrm{final}}$ & avg & BWT \\
\midrule
$T{=}10$ & SI ($100$) & 40.01 & 13.23 & 39.26 & $-0.83$ \\
$T{=}10$ & MAS ($1$)   & 85.71 & 87.00 & 66.61 & $-21.21$ \\
$T{=}20$ & SI ($100$) & 25.80 & 11.19 & 25.14 & $-0.69$ \\
$T{=}20$ & MAS ($1$)   & 86.69 & 87.91 & 55.75 & $-32.56$ \\
\bottomrule
\end{tabular}
\end{table}

\subsection{The flatness survives anchor budgets}
\label{sec:budget}

Because segment-resolved anchors might pay off only when the anchor
\emph{stack} is bounded --- fewer anchors plausibly sharpen the value of
each --- we crossed granularity with FIFO budgets $M \in \{2, 5\}$ at
$\lambda{=}1000$ (preregistered: alive iff $\ge 4/6$ cells significant
positive and none significantly worse on BWT). The result is the same null,
tighter: across all eight $\{$dataset$\}\times\{M\}\times\{T\}$ cells the
branch$-$neuron $\Delta\ell$ lies within $\pm 0.44$ with no cell significant
positive (one cell, Fashion $M{=}2$, $T{=}10$, is significantly
\emph{negative} at $-0.33$, $t(9){=}-2.3$; Table~\ref{tab:budget}).
Budgets shift both granularities identically, consistent with the shared
mechanism analysis below.

\begin{table}[t]
\centering\small
\caption{Granularity under FIFO anchor budgets ($\lambda{=}1000$, $n{=}10$
paired). No branch$-$neuron difference is significant positive in any of
the eight cells.}
\label{tab:budget}
\begin{tabular}{llrrrr}
\toprule
cell & $M$ & neuron $\ell$ & branch $\ell$ & $\Delta\ell$ (t) & $\Delta$BWT (t) \\
\midrule
MNIST $T{=}10$  & 2 & 83.06 & 83.01 & $-0.05\;(-1.0)$ & $+0.22\;(+1.1)$ \\
MNIST $T{=}10$  & 5 & 77.09 & 77.02 & $-0.07\;(-0.5)$ & $+0.15\;(+1.0)$ \\
MNIST $T{=}20$  & 2 & 80.23 & 80.32 & $+0.09\;(+1.0)$ & $+0.39\;(+1.9)$ \\
MNIST $T{=}20$  & 5 & 68.30 & 68.19 & $-0.11\;(-0.8)$ & $+0.28\;(+1.5)$ \\
Fashion $T{=}10$ & 2 & 66.66 & 66.33 & $-0.33\;(-2.3)^{*}$ & $-0.41\;(-1.2)$ \\
Fashion $T{=}10$ & 5 & 58.90 & 58.74 & $-0.16\;(-0.9)$ & $+0.43\;(+1.2)$ \\
Fashion $T{=}20$ & 2 & 62.61 & 62.72 & $+0.10\;(+0.4)$ & $-0.18\;(-0.3)$ \\
Fashion $T{=}20$ & 5 & 49.79 & 50.23 & $+0.44\;(+1.5)$ & $-0.33\;(-0.7)$ \\
\bottomrule
\end{tabular}
\end{table}

% =====================================================================


\section{Why the ladder is flat: auditing segment--task specialization}
\label{sec:audit}

The null in Sec.~\ref{sec:granularity} would be easy to misread as ``the
segment machinery is redundant.'' The audit says something stronger: the
specialization that segment-resolved protection would need does not exist,
and part of the evidence that appeared to show it was an artifact.

\subsection{A measurement artifact: dead neurons fabricate dominance}
\label{sec:artifact}

Our first-pass instrumentation recorded, for every neuron, the argmax over
its segment activations under each task's data. It showed segment 0
``dominating'' $58$--$63\%$ of neurons --- a striking positional bias that
initially looked like the failure of task-specific routing. The correction:
under the production sampling budget ($500$ examples), $56$--$63\%$ of
neurons are \emph{never activated}; for a dead neuron all segment
activations are exactly zero, argmax breaks ties toward index $0$, and the
neuron is silently booked to segment 0 (Fig.~\ref{fig:artifact}).
Filtering to active neurons collapses the apparent bias to chance
($\approx 0.10 = 1/T$ at $T{=}10$; untrained networks show no positional
preference either, $0.096$--$0.111$ over three seeds). The death itself is
real and grows with the sequence: the inactive fraction jumps from $0.02$
to $0.35$ within the first phase of training and stabilizes near $0.6$
(Fig.~\ref{fig:deadfrac}) --- under kWTA, neurons pushed out of the
winner set early are never re-selected by any later task. We flag this as
a methodological warning for the active-dendrite literature: dominance
statistics computed over all neurons conflate routing with mortality, and
at production sampling budgets the conflation is an order of magnitude
larger than the signal.

\begin{figure}[t]
\centering
\includegraphics[width=.55\linewidth]{fig_artifact.pdf}
\caption{The dead-neuron artifact. Naive argmax over \emph{all} neurons
reports segment-0 dominance of $\approx 60\%$ after training; the same
statistic over active neurons only, and both statistics on untrained
networks, sit at chance ($1/T$).}
\label{fig:artifact}
\end{figure}

\begin{figure}[t]
\centering
\includegraphics[width=.5\linewidth]{fig_deadfrac.pdf}
\caption{Inactive-neuron fraction across the task sequence (end-of-phase
means). Neurons lost in early phases are never recovered; the active
substrate available to routing is roughly halved by phase 2 and
stabilizes near $40\%$.}
\label{fig:deadfrac}
\end{figure}

\subsection{Anti-specialization among active neurons}
\label{sec:anti}

With the artifact controlled, the true routing picture is the opposite of
specialization. Define $\own$ as the probability that a neuron's dominant
segment is the currently measured task's own segment, over active neurons.
At the \emph{start} of each phase, $\own$ sits at chance ($1/T$): nothing
in the context-gated routing maps tasks to segments. Training within the
phase then pushes it \emph{below} chance --- by the end of phase $p$, $\own$
drops to $0.058$--$0.085$ at $T{=}10$ (chance $0.10$) and, late in a
$T{=}20$ sequence, to $0.005$--$0.02$ against a chance of $0.05$, i.e.\ five
to ten times \emph{below} random assignment (Fig.~\ref{fig:evolution}).
The own-segment gating block does train --- it moves by $6$--$16\%$ of its
norm per phase --- yet this movement fails to convert into dominance: win
rates still end each phase below chance. Segment--task alignment is not
merely absent; its absence persists through active training of the very
parameters that would encode it.

\begin{figure}[t]
\centering
\includegraphics[width=\linewidth]{fig_evolution.pdf}
\caption{Anti-specialization over training. Within each phase, the
own-segment win rate over active neurons starts at chance (dashed) and
ends below it; the deficit deepens with sequence position (rightmost
points of each curve), most severely at $T{=}20$.}
\label{fig:evolution}
\end{figure}

\subsection{Fisher mass is diffuse across segments}

Segment-resolved anchors also lack a substrate to hold onto: for trained
anchors, the fraction of a neuron's Fisher mass attributable to its top
segment is only $11$--$18\%$ (mean over anchor--neuron pairs). Even a
perfectly attributed segment-level penalty would spread nearly all of its
protection across segments the anchor task never drove. This measurement
closes the mechanistic account of the Sec.~\ref{sec:granularity} null from
two sides: attribution is self-limiting (an anchor's own penalty pins the
shared weights it measures, so its tracked references barely move), and
the per-segment mass that attribution would partition is diffuse to begin
with.

% =====================================================================


\section{Causal probe: restoring own-segment advantage}
\label{sec:causal}

\paragraph{How to read this section.} The intervention studied here
produces a large, consistent, saturating effect on the learning term, and
in the version of this work that first reported it, that effect was taken
to be a causal repair of the missing specialization. Two later sections
of this manuscript do not sustain that reading: on a convolutional
active-dendrite backbone the dose--response collapses to noise
(Sec.~\ref{sec:conv}), and a mechanism probe shows that the repaired gate
is \emph{context-degenerate} --- the bias works by saturating a gate that
was never context-dependent, not by restoring routing
(Secs.~\ref{sec:mech}--\ref{sec:ablation}). We report the dose--response
in full anyway, for a reason that only becomes clear later: its
\emph{shape} is what rules out the dying-gate explanation. A repair that
revived dying gates would be expected to weaken when mortality changes;
this one does not, which points elsewhere. A reader who wants only the
outcome may skip to Sec.~\ref{sec:conv}; a reader assessing the mechanism
should read the dose--response first.

\subsection{A one-parameter intervention}

If the missing specialization is load-bearing, causally restoring it should
improve the trade-off in Eq.~\eqref{eq:identity}. We add a constant bias $b$
to the own-segment dendritic input: $u_{j,t} \mathrel{+}= b$ whenever the
context selects task $t$. The bias is active in training, evaluation, and
anchor construction (the Fisher measurement that drives usage-localized
protection); the dominant-segment win-rate instrument deliberately reads
the raw, bias-free gates, so that every routing statistic in this paper
describes what the network itself encodes rather than what the bias
forces (Sec.~\ref{sec:compensation}). The intervention is deliberately
dumb --- no learned parameters, no schedule, no per-neuron structure ---
so that any effect is attributable to the routing prior alone. Doses
$b \in \{0, 1, 3, 10\}$; all comparisons paired against $b{=}0$ on
identical seeds.

\subsection{Result: the learning term and accuracy improve everywhere}
\label{sec:dosestat}

\paragraph{Statistics.} All tests in this section are paired $t$-tests over $n=10$ seeds sharing permutation streams and initialization. Because the dose grid spans three metrics across $22$ cells, we report Holm-corrected $p$-values within each metric family ($\Delta\ell$, $\Delta\mathrm{avg}$, $\Delta\mathrm{BWT}$) in addition to the raw values; effect sizes are paired standardized means ($d_z=\bar{d}/s_d$) and intervals are $95\%$ CIs for the paired mean difference. The correction removes two cells from the $\Delta\ell$ family and four from the $\Delta\mathrm{avg}$ family (both on Fashion-MNIST at $T=10$, where the effect is smallest); it does not touch the qualitative result, which rests on the MNIST cells and on the $T=20$ monotone dose--response.

Tables~\ref{tab:dose10} and~\ref{tab:dose20} give the dose--response;
Fig.~\ref{fig:dose} plots it. The full grid is
$2$ datasets $\times$ $2$ depths $\times$ $2$ granularities $\times$
$3$ doses $=24$ cells; $22$ of them are reported here (the two
$T{=}20$ branch cells at $b{=}1$ were not collected), and across the $22$
cells we do report, $\Delta\ell$ is positive in every cell and
significant in
$16/22$; $\ell$ and final accuracy are simultaneously
significant in $16/22$, and every grid cell has at least one dose
significant on both terms. The effect grows with depth: at $T{=}20$ MNIST,
$b{=}10$ lifts $\ell$ by $+4.56$ ($t(9){=}8.1$) and final accuracy by
$+3.81$ ($t(9){=}10.6$) over the already usage-localized baseline; the
$T{=}20$ learning curve is monotone in dose
($\Delta\ell$ $+2.14$ / $+3.72$ / $+4.56$ for $b{=}1/3/10$, all
$t \ge 3.9$), and $b{=}3$ alone already captures $81\%$ of the
$b{=}10$ effect there ($81$--$113\%$ across the six dataset/depth/granularity
arms we ran at all three doses). Forgetting
improves broadly at $T{=}10$ ($\Delta\mathrm{BWT}$ up to $+1.65$;
significantly positive in $6/12$ cells --- all six MNIST cells --- with
      no Fashion cell surviving correction), and
partially trades back at $T{=}20$, where $3/10$ cells show a significant
mild BWT cost ($-0.66$ to $-0.95$) and the remaining seven are
statistically indistinguishable from baseline.

Two robustness notes. First, the $22/22$ sign uniformity of $\Delta\ell$
is itself the strongest guard against multiplicity: a sign test gives
$p = 2^{-22}$, orders of magnitude below any family-wise correction over
the grid. Second, the intervention also reduces neuron mortality ---
measured on the raw, bias-free network, the mean dead fraction drops by
$3.9$pp (MNIST) and $2.3$pp (Fashion) at $T{=}10$ --- and per-run
$\Delta\ell$ correlates with the mortality reduction ($r{=}-0.49$ and
$-0.68$; both correlations are computed on the ten paired seeds of the
$T{=}10$ neuron arm at $b{=}10$, the same cell in which the dead-fraction
shift is quoted). We read this as one causal chain rather than a confound:
raising the own-segment gate raises $\mathrm{sigmoid}(u)$ during
training, which increases kWTA survival, so forcing the routing prior is
also what keeps neurons alive; the routing and capacity channels are two
faces of a single intervention, not competing explanations. That chain is
a \emph{within-configuration} one, and it does not extend between
configurations: across the eight cells of the width$\times$normalization
scan the mortality of a configuration and the size of its repair gain are
uncorrelated, and in the two cells where mortality is essentially
eliminated the gain is the smallest and can be negative
(Sec.~\ref{sec:zeromort}, Table~\ref{tab:scanmort}).

\begin{table}[t]
\centering\small
\caption{Causal probe, $T={10}$, $\lambda=1000$: paired differences vs.\ $b=0$ ($n=10$ paired seeds). $\dagger$: Holm-corrected $p<0.05$ within the $\Delta\ell$ family; $\ddagger$: within the $\Delta\mathrm{avg}$ family; $\S$: within the $\Delta\mathrm{BWT}$ family. Intervals are 95\% CI for the paired mean difference; $d_z$ is the paired standardized effect size.}
\label{tab:dose10}
\resizebox{\textwidth}{!}{%
\begin{tabular}{llrrrrrl}
\toprule
dataset & arm ($b$) & $\Delta\ell$ & $\Delta\mathrm{avg}$ & $\Delta\mathrm{BWT}$ & $p_{\mathrm{Holm}}(\Delta\ell)$ & $d_z$ & 95\% CI \\
\midrule
fashion & branch ($1$) & $+0.94{}$ & $+2.08{}$ & $+1.26{}$ & $0.247$ & $+0.67$ & $[-0.06, +1.94]$ \\
fashion & branch ($3$) & $+0.81{}$ & $+1.79{}$ & $+1.09{}$ & $0.248$ & $+0.54$ & $[-0.27, +1.88]$ \\
fashion & branch ($10$) & $+0.97{}$ & $+1.03{}$ & $+0.06{}$ & $0.247$ & $+0.64$ & $[-0.11, +2.06]$ \\
fashion & neuron ($1$) & $+0.87{}$ & $+1.03{}$ & $+0.18{}$ & $0.248$ & $+0.53$ & $[-0.31, +2.06]$ \\
fashion & neuron ($3$) & $+1.23{}$ & $+2.56{}$ & $+1.48{}$ & $0.140$ & $+0.86$ & $[+0.21, +2.25]$ \\
fashion & neuron ($10$) & $+1.39{}$ & $+0.69{}$ & $-0.77{}$ & $0.192$ & $+0.77$ & $[+0.09, +2.68]$ \\
mnist & branch ($1$) & $+2.18{\dagger}$ & $+3.29{\ddagger}$ & $+1.23{\S}$ & $0.000$ & $+2.48$ & $[+1.55, +2.81]$ \\
mnist & branch ($3$) & $+2.90{\dagger}$ & $+4.30{\ddagger}$ & $+1.57{\S}$ & $0.000$ & $+3.73$ & $[+2.34, +3.45]$ \\
mnist & branch ($10$) & $+3.16{\dagger}$ & $+4.17{\ddagger}$ & $+1.13{\S}$ & $0.000$ & $+3.33$ & $[+2.48, +3.84]$ \\
mnist & neuron ($1$) & $+1.99{\dagger}$ & $+3.40{\ddagger}$ & $+1.56{\S}$ & $0.000$ & $+2.56$ & $[+1.44, +2.55]$ \\
mnist & neuron ($3$) & $+2.68{\dagger}$ & $+4.08{\ddagger}$ & $+1.56{\S}$ & $0.000$ & $+4.21$ & $[+2.22, +3.13]$ \\
mnist & neuron ($10$) & $+3.32{\dagger}$ & $+4.80{\ddagger}$ & $+1.65{\S}$ & $0.000$ & $+4.64$ & $[+2.81, +3.83]$ \\
\bottomrule
\end{tabular}
}
\end{table}


\begin{table}[t]
\centering\small
\caption{Causal probe, $T={20}$, $\lambda=1000$: paired differences vs.\ $b=0$ ($n=10$ paired seeds). $\dagger$: Holm-corrected $p<0.05$ within the $\Delta\ell$ family; $\ddagger$: within the $\Delta\mathrm{avg}$ family; $\S$: within the $\Delta\mathrm{BWT}$ family. Intervals are 95\% CI for the paired mean difference; $d_z$ is the paired standardized effect size.}
\label{tab:dose20}
\resizebox{\textwidth}{!}{%
\begin{tabular}{llrrrrrl}
\toprule
dataset & arm ($b$) & $\Delta\ell$ & $\Delta\mathrm{avg}$ & $\Delta\mathrm{BWT}$ & $p_{\mathrm{Holm}}(\Delta\ell)$ & $d_z$ & 95\% CI \\
\midrule
fashion & branch ($3$) & $+1.44{\dagger}$ & $+1.98{\ddagger}$ & $+0.57{}$ & $0.040$ & $+1.14$ & $[+0.54, +2.34]$ \\
fashion & branch ($10$) & $+2.18{\dagger}$ & $+1.96{\ddagger}$ & $-0.24{}$ & $0.004$ & $+1.73$ & $[+1.28, +3.09]$ \\
fashion & neuron ($1$) & $+1.78{\dagger}$ & $+1.58{\ddagger}$ & $-0.21{}$ & $0.025$ & $+1.29$ & $[+0.79, +2.77]$ \\
fashion & neuron ($3$) & $+2.19{\dagger}$ & $+2.09{\ddagger}$ & $-0.11{}$ & $0.003$ & $+1.84$ & $[+1.34, +3.04]$ \\
fashion & neuron ($10$) & $+1.94{\dagger}$ & $+1.03{\ddagger}$ & $-0.95{}$ & $0.004$ & $+1.72$ & $[+1.13, +2.75]$ \\
mnist & branch ($3$) & $+3.66{\dagger}$ & $+3.03{\ddagger}$ & $-0.66{}$ & $0.000$ & $+2.95$ & $[+2.77, +4.55]$ \\
mnist & branch ($10$) & $+4.13{\dagger}$ & $+3.48{\ddagger}$ & $-0.68{}$ & $0.000$ & $+2.46$ & $[+2.93, +5.33]$ \\
mnist & neuron ($1$) & $+2.14{\dagger}$ & $+2.78{\ddagger}$ & $+0.68{}$ & $0.031$ & $+1.22$ & $[+0.89, +3.39]$ \\
mnist & neuron ($3$) & $+3.72{\dagger}$ & $+3.25{\ddagger}$ & $-0.49{}$ & $0.000$ & $+2.36$ & $[+2.59, +4.84]$ \\
mnist & neuron ($10$) & $+4.56{\dagger}$ & $+3.81{\ddagger}$ & $-0.79{}$ & $0.000$ & $+2.57$ & $[+3.29, +5.83]$ \\
\bottomrule
\end{tabular}
}
\end{table}


\begin{figure}[t]
\centering
\includegraphics[width=\linewidth]{fig_dose.pdf}
\caption{Dose--response of the own-segment bias (neuron-resolved arm,
paired means vs.\ $b{=}0$). The learning term (left) and final accuracy
(middle) improve at every dose, dataset, and depth; BWT (right) improves
at $T{=}10$ and is mildly negative at $T{=}20$ on MNIST. Note the
saturating shape: $b{=}3$ alone already captures most of the $b{=}10$
effect ($81$--$113\%$ across datasets and depths).}
\label{fig:dose}
\end{figure}



\subsection{What the response is, and is not, made of}
\label{sec:compact}
\label{sec:compensation}
\label{sec:taskrel}
\label{sec:sigdrift}

Three further measurements characterize the intervention, and we report
them compactly because, as the forward reference above notes, their
role is diagnostic rather than confirmatory: they are what later rules
out the dying-gate explanation.

\textbf{The network absorbs the intervention rather than opposing it.}
Under per-phase signed instrumentation of the context-$t$ gating column
($n=10$ seeds, $T{=}10$ MNIST), we saw the own-segment entries drift
\emph{positive} in all ten phases under $b{=}10$ (one-sample $t$ over
seeds: $170$, $54$, $25$, $22$, $14$, $16$, $8.0$, $14$, $8.1$, $14$;
per-phase mean displacement decaying from $+0.027$ to $+0.004$). Under
$b{=}0$ the same quantity is small but not absent: its per-phase means
peak at $+0.0041$, a factor of $6.5$ below the $b{=}10$ peak, and it is
nominally significant in phases $0$--$7$ at $t$ between $2.8$ and $23$
before turning negative in phase $8$. We therefore compare magnitudes
rather than significance thresholds here: the network drifts \emph{with}
the bias at every dose and by an amount that scales with it, which is
the opposite of what adversarial compensation would produce.
The raw own gate stays within $|u| \le 0.325$ across all seeds and phases
--- thirty times below the imposed $b{=}10$ and three times below even
$b{=}1$ --- so the offset remains fully in effect throughout training.
Adversarial compensation predicts the opposite displacement; it does not
occur.

An earlier version of this paragraph quoted the same per-phase
$t$ statistics over three seeds. That statistic is not reproducible: the
per-seed means agree across machines to $10^{-4}$, but a $t$ over three
numbers is dominated by the sample standard deviation, and float-level
noise moved the values by several units between runs. We report the
ten-seed version instead, and the qualitative reading is unchanged.

\textbf{The saturation is a scale effect, not a partial restoration.} The
own-segment win rate saturates at $13$--$19\%$ regardless of dose,
because any effective dose dwarfs the raw gate scale ($|u|\le 0.32$): the
bias forces the same training-time routing at every dose tested, so only
a minority of neuron-tasks flip to the own-segment pathway --- and that
partial restoration already yields most of the gain. This matters for
the mechanism: the effect is not proportional to how much routing is
restored.

\textbf{The benefit requires learnable structure.} In a task-relevance
probe (ten interleaved phases: two in-domain tasks, two rotated-count
tasks, six zero-relevance tasks; $b\in\{0,1,3,10\}$, $n=10$ paired seeds),
zero-relevance tasks sit at chance at every dose (per-task accuracy
10.2--11.4 across the whole
sweep), the gap between the high-relevance group
($\ell\approx94$) and the medium group
($\ell\approx91$) narrows monotonically with
dose ($+3.0$ at $b{=}0$ to $+1.6$ at
$b{=}10$), and the high-relevance group still shows the largest absolute
gain. The bias therefore behaves like a structured capacity channel whose
marginal benefit is largest where the network's own routing is weakest, and
delivers nothing where there is nothing to amplify.

Two caveats attach to this probe, and we state them because the earlier
version of this line of work did not. First, the zero-relevance group is
built from random labels, whose accuracy is chance by construction; the
observation that the bias does not move it is therefore weaker evidence
than it looks. We therefore added a non-trivial zero-relevance condition:
half of the zero-relevance tasks use labels rotated \emph{within} each
class across tasks, so the class structure remains learnable and only the
task--label mapping changes. Second, the probe was re-run for this
manuscript because the original run's script and per-seed records were not
in the artifact, so the numbers here come from a fresh run rather than from
the one quoted in the earlier version of this work; the qualitative
pattern (monotone narrowing of the gap, chance-level performance on the
zero-relevance group) reproduces, the absolute values do not, and we report
the reproduced ones.

These three measurements together --- absorption, dose-invariant
saturation, and structure-dependence --- describe a mechanism that is not
restoring segment--task routing, and they are the first hint of the
account given in Secs.~\ref{sec:mech}--\ref{sec:ablation}. Full detail,
including the per-phase drift tables and the per-group accuracy matrices,
is in the supplementary material.

\subsection{Why forgetting trades back at depth}
\label{sec:depthbwt}

At $T{=}10$ the bias reduces forgetting; at $T{=}20$ MNIST it mildly
increases it. The probe's measurement explains the sign flip. With the
bias active, a fixed minority of neurons routes through the own segment in
\emph{every} phase; over a longer sequence, later phases increasingly
share this minority substrate, so late-task updates overlap earlier
anchors on the same (now doubly used) neurons --- $\ell$ rises because
more neurons are available early, but the shared minority pays for it at
depth. The saturation of $\own$ (Sec.~\ref{sec:compensation}) caps how
much substrate the repair can recruit; at $T{=}20$ the cap binds.

% =====================================================================



% ===== S3: 架构边界 =====


\section{The repair does not transfer: a convolutional active-dendrite backbone}
\label{sec:conv}

\subsection{Architecture and equivalence of the instrument stack}

ConvADNet replaces only the feature path of the bare MLP's ADNet:
\[
x \mapsto [\mathrm{Conv}{\text -}\mathrm{GN}{\text -}\mathrm{ReLU}]
\times 4 \mapsto \mathrm{GAP} \mapsto t \in \mathbb{R}^{2048},
\]
with strides $(1,2,2,2)$ for $28{\times}28$ inputs. The dendritic
gate, kWTA ($k{=}64$), readout, protected gradients, anchor
construction (Fisher $\times$ usage), evaluation protocol, and the
raw-gate routing instrument are \emph{byte-identical} to the bare-MLP arm: the gate
reads $\mathrm{dend}(\mathrm{ctx})$ directly, so every routing
statistic measures what the network encodes, not what a bias forces.
Hyperparameters are unchanged (Adam $10^{-3}$, $15$ epochs,
$\lambda{=}1000$, seeds $0$--$9$ per cell). $160$ runs total
($b \in \{0,1,3,10\}$ $\times$ $2$ datasets $\times$ $T \in
\{10,20\}$); zero failures.

\subsection{G1 fails: the dose--response collapses}
\label{sec:convstat}

\paragraph{Statistics.} Paired $t$-tests over $n=10$ seeds, with Holm correction applied within each metric family across the $12$ conv dose cells. Uncorrected, $2/12$ cells reach $p<0.05$ on $\Delta\ell$; after correction \emph{none} does. We report the corrected result as the primary one, because the pre-registered criterion was a count over cells and a count over uncorrected tests is the wrong statistic for it. The conclusion (G1 fails) is unchanged either way, and is in fact strengthened: the conv dose--response is not a weak effect that shrinks under correction but an effect that is not detectable at this sample size, with 95\% CIs that in two
cells admit effects up to $+3$ points.

Table~\ref{tab:convdose} reports paired differences against $b{=}0$.
Against the pre-registered bar ($\ge 15/20$ cells significant), the
result is unambiguous: \textbf{$2/12$ cells} reach $p<0.05$ on $\Delta\ell$
uncorrected and \textbf{$0/12$ after Holm correction}, largest effect
$+1.46$ ($t{=}2.1$, uncorrected), sign-consistency $11/12$. The third
pre-registered G1 condition --- $\Delta\ell$ and $\Delta\mathrm{avg}$
simultaneously significant --- fails most completely of the three:
\textbf{$0/12$ cells}, i.e.\ no conv dose cell moves the learning term
and the final accuracy in the same significant direction at once. In the
five conv cells with a matched bare-MLP counterpart, the MLP's effects were
$4$--$56\times$ larger with paired $|t|$ up to $14.7$
(Table~\ref{tab:three}, column $\Delta\ell$). Worse, final accuracy
\emph{reverses}: at fashion $T{=}10$, $b{=}10$ costs
$\Delta\mathrm{avg} = -2.29$ ($t{=}-3.1$) --- the repair buys a
non-significant learning sliver at a significant accuracy cost.

\begin{figure}[t]
\centering
\includegraphics[width=\linewidth]{fig_dose3.pdf}
\caption{Dose--response of the own-segment bias across the three
configurations ($T{=}10$, paired means vs.\ each configuration's own
$b{=}0$). The bare MLP (blue) shows the bare-MLP dose--response; GN-MLP
(orange) collapses to a fifth of it; ConvADNet (green) to noise.}
\label{fig:dose3}
\end{figure}

\begin{table}[t]
\centering\small
\caption{ConvADNet dose--response, paired vs.\ $b=0$ ($n=10$; Holm-corrected $p$ within each metric family; 95\% CI for the paired mean difference; $d_z$ the paired standardized effect size). No cell survives correction on any metric, which is the pre-registered G1 criterion evaluated conservatively.}
\label{tab:convdose}
\resizebox{\textwidth}{!}{%
\begin{tabular}{llrrrrrl}
\toprule
cell & arm ($b$) & $\Delta\ell$ & $\Delta\mathrm{avg}$ & $\Delta\mathrm{BWT}$ & $p_{\mathrm{Holm}}(\Delta\ell)$ & $d_z$ & 95\% CI \\
\midrule
fashion $T=10$ & $1$ & $-0.06$ & $+0.22$ & $+0.31$ & $1.000$ & $-0.22$ & $[-0.25, +0.14]$ \\
fashion $T=20$ & $1$ & $+0.31$ & $+0.66$ & $+0.36$ & $0.081$ & $+1.11$ & $[+0.11, +0.51]$ \\
mnist $T=10$ & $1$ & $+0.17$ & $-1.64$ & $-2.02$ & $1.000$ & $+0.12$ & $[-0.87, +1.22]$ \\
mnist $T=20$ & $1$ & $+0.63$ & $-0.24$ & $-0.92$ & $1.000$ & $+0.24$ & $[-1.21, +2.48]$ \\
fashion $T=10$ & $3$ & $+0.02$ & $-0.89$ & $-1.01$ & $1.000$ & $+0.03$ & $[-0.48, +0.52]$ \\
fashion $T=20$ & $3$ & $+0.29$ & $-0.51$ & $-0.84$ & $0.540$ & $+0.72$ & $[+0.00, +0.58]$ \\
mnist $T=10$ & $3$ & $+0.30$ & $-0.99$ & $-1.43$ & $1.000$ & $+0.18$ & $[-0.87, +1.47]$ \\
mnist $T=20$ & $3$ & $+1.46$ & $-0.26$ & $-1.81$ & $0.642$ & $+0.67$ & $[-0.11, +3.02]$ \\
fashion $T=10$ & $10$ & $+0.36$ & $-2.29$ & $-2.95$ & $1.000$ & $+0.45$ & $[-0.21, +0.94]$ \\
fashion $T=20$ & $10$ & $+0.31$ & $-0.87$ & $-1.24$ & $0.642$ & $+0.66$ & $[-0.03, +0.64]$ \\
mnist $T=10$ & $10$ & $+0.07$ & $-0.84$ & $-1.01$ & $1.000$ & $+0.04$ & $[-1.23, +1.38]$ \\
mnist $T=20$ & $10$ & $+1.28$ & $+0.30$ & $-1.04$ & $0.682$ & $+0.61$ & $[-0.22, +2.79]$ \\
\bottomrule
\end{tabular}
}
\end{table}


\subsection{G2 fails with a clean cause: zero mortality}
\label{sec:zeromort}

The dominant-segment instrument reads identically zero artifact:
naive share $=$ active share $=$ chance $= 1/T$ in every cell
(Table~\ref{tab:three}, column ``dead''). The cause is visible in the
per-phase instrumentation: ConvADNet's dead-gate fraction is
\emph{exactly $0.000$ in every phase of every run}, against the bare
MLP's $0.572$ . GroupNorm's affine output keeps
$|u| > 0$; there are no argmax ties over dead gates because there are
no dead gates (Fig.~\ref{fig:mortality}).

The same conclusion holds at the level of configurations rather than
runs, which is where the dying-gate hypothesis was proposed. We
re-ran the full width$\times$normalization factorial scan with a
mortality column added ($160$ runs: $2$ datasets $\times$ $2$ widths
$\times$ $2$ normalization settings $\times$ $2$ doses $\times$ $10$
seeds; Table~\ref{tab:scanmort}). The repair gain in final accuracy
and the mortality of the repaired arm are, across the eight cells,
uncorrelated: Pearson $r = -0.11$ ($p = 0.79$), Spearman
$\rho = +0.05$ ($p = 0.91$). More than a flat correlation, the ordering
inverts locally. The two GroupNorm cells have the \emph{lowest}
mortality of the four configurations on their dataset ($0.000$ and
$0.020$ on MNIST and Fashion-MNIST, against $0.535$--$0.948$ without
it) and the \emph{smallest} repair gain, which is negative on both
datasets ($-2.74$ and $-6.64$ points); the highest-mortality cells are
the ones with the largest positive gain ($+4.62$ and $+2.13$). A
dying-gate explanation predicts that removing dead gates should
restore the repair's benefit; in the cell where mortality is
essentially eliminated, the benefit is at its smallest and can be
negative.

This is consistent with, but weaker than, the within-cell correlation
reported above: relieving mortality on a given run helps that run's
learning term, yet the configuration with the most dead neurons is not
the configuration with the biggest repair. The causal chain we
describe is therefore a within-configuration one; it does not extend to
the between-configuration comparison, and the dying-gate hypothesis is
not supported at that level either.

A measurement caveat we state explicitly: the two configurations
report mortality under \emph{different} zero-conditions. The MLP
instrument (inherited unchanged from the bare-MLP arm) marks a neuron dead if its
post-kWTA activity is never positive across the usage window ---
mortality there means ``never wins the top-$k$ competition''. The
conv instrument marks a neuron dead if its gate row is identically
zero. Both are the natural operationalizations in their families,
but they are not the same statistic, and the conv family's
post-kWTA ``never wins'' rate was not measured. What the comparison
supports is therefore the narrower claim we make: the bare-MLP
\emph{artifact mechanism} (argmax ties over exactly-zero gates) has
no substrate on the conv configuration --- its precondition
(numerically zero gates) never occurs. The artifact is thus
\textbf{architecture-bound, not instrument-bound}: the same reading
instrument that the bare-MLP arm showed contaminated reads clean here, for reasons
that have nothing to do with the instrument.

\begin{figure}[t]
\centering
\includegraphics[width=0.72\linewidth]{fig_mortality.pdf}
\caption{Dead-gate fraction at $b{=}0$ across configurations and
datasets ($T{=}10$, mean over phases and seeds). Conv+GroupNorm
mortality is exactly zero; GroupNorm inserted into the MLP
\emph{raises} mortality. Mortality is a property of the architecture,
not of the intervention.}
\label{fig:mortality}
\end{figure}

\subsection{The conv baseline itself}

ConvADNet's unbiased baseline learns more and forgets more than the
MLP's: $\ell_0$ $84.1$/$75.3$ (vs.\ $74.3$/$56.8$) but
$\mathrm{BWT}_0$ $-50.4$/$-32.8$ (vs.\ $-11.6$/$-16.2$) at $T{=}10$.
The protection strength $\lambda{=}1000$, calibrated for the MLP's
parameter scale, does not carry over --- itself a transfer finding we
report rather than tune away (see Sec.~\ref{sec:boundary}).

% =====================================================================


\section{The dying-gate hypothesis, tested and rejected}
\label{sec:gnmlp}

The natural reading of Sec.~\ref{sec:conv} is: the repair worked in
the bare-MLP arm because it revived dying gates, and conv gates do not die. This
``dying-gate hypothesis'' makes a sharp same-architecture prediction:
insert a normalization layer into the MLP (features $t =
\mathrm{GN}(\mathrm{ff}(x))$; $80$ runs,
seeds $0$--$9$, $b \in \{0,1,3,10\}$, both datasets, $T{=}10$) ---
mortality should fall and the repair should weaken accordingly.
One deviation from strict matching is disclosed rather than hidden:
the GN-MLP uses $2000$ hidden units where the bare MLP uses $256$, so
the control differs in width as well as normalization; a factor scan
below separates the two.

The prediction fails in both halves (Table~\ref{tab:three}, middle
rows). Mortality \emph{rises} ($0.765$/$0.855$ vs.\ bare's
$0.572$/$0.603$), for a reason the architecture makes plain:
GroupNorm rescales the features $t$ symmetrically about zero, and the
kWTA competition keeps the $64$ \emph{largest} values of $t \odot
\sigma(d)$ --- so every neuron whose normalized feature is
predominantly negative can never enter the top-$k$ and counts as
dead under the usage instrument. GroupNorm does not touch the gate
zero-set $u$; it manufactures post-kWTA death through the feature
sign structure instead. Meanwhile the repair gain \emph{collapses
anyway}: to
$+0.57$ (retention $17\%$ of the bare-MLP gain) on MNIST and $+0.36$
($26\%$) on Fashion. And final accuracy reverses on both datasets
($\Delta\mathrm{avg} = -3.88$, $t{=}-5.3$; $-6.55$, $t{=}-4.4$).

Mortality and gain therefore do not track each other: one configuration
(GN-MLP) has the highest mortality and the least gain; another
(ConvADNet) has zero mortality and no gain; the bare MLP has middle
mortality and full gain. Whatever the repair feeds on, it is not gate
death per se.

\begin{table}[t]
\centering\small
\caption{The three configurations at $T{=}10$ ($n{=}10$; paired
vs.\ each configuration's own $b{=}0$). ``dead'' = mean dead-gate
fraction; ``HR$_\ell$'' = $100 - \ell_0$. The repair's
$\Delta\ell$ tracks baseline weakness, not mortality.}
\label{tab:three}
\begin{tabular}{llrrrrr}
\toprule
dataset & config & $\ell_0$ & HR$_\ell$ & dead &
$\Delta\ell_{b=10}$ & $\Delta\mathrm{avg}_{b=10}$ \\
\midrule
MNIST & bare-MLP & $74.3$ & $25.7$ & $0.572$ &
$+3.32^{*}$ & $+4.80^{*}$ \\
MNIST & GN-MLP & $93.5$ & $6.5$ & $\mathbf{0.765}$ &
$+0.57^{*}$ & $-3.88^{*}$ \\
MNIST & ConvADNet & $84.1$ & $15.9$ & $\mathbf{0.000}$ &
$+0.07$ & $-0.84$ \\
\midrule
Fashion & bare-MLP & $56.8$ & $43.2$ & $0.603$ &
$+1.39^{*}$ & $+0.69$ \\
Fashion & GN-MLP & $79.3$ & $20.7$ & $\mathbf{0.855}$ &
$+0.36^{*}$ & $-6.55^{*}$ \\
Fashion & ConvADNet & $75.3$ & $24.7$ & $\mathbf{0.000}$ &
$+0.36$ & $-2.29^{*}$ \\
\bottomrule
\end{tabular}
\end{table}

\paragraph{Width vs.\ normalization: a $2{\times}2$ factor scan.}
The width deviation disclosed above is not merely acknowledged --- it
is measured. We re-ran the bare-MLP vs.\ GN-MLP comparison as a full
$2{\times}2$ factor design: hidden width $n_h \in \{256, 2000\}$
crossed with normalization $\{$none, $\mathrm{GN}\}$, at $b \in
\{0,1,3,10\}$, both datasets, $10$ seeds --- $320$ runs. The scan is
an independent reimplementation, and its GN-MLP cells replicate the
original $80$-run records: $\ell_0$ $93.4$/$79.2$ vs.\
$93.5$/$79.3$, gain $+0.58$/$+0.49$ vs.\ $+0.57$/$+0.36$.

\begin{table}[t]
\centering\small
\caption{Width $\times$ normalization factor scan ($n{=}10$; paired
vs.\ each cell's own $b{=}0$; $b{=}10$ learning-term gain with paired
$t$). The GN collapse appears at both widths; the width effect on the
plain trunk has the opposite sign.}
\label{tab:twobytwo}
\begin{tabular}{llrr}
\toprule
dataset & config & $\ell_0$ & $\Delta\ell_{b=10}$ \\
\midrule
MNIST & $n_h{=}256$, plain & $74.3$ & $+3.32^{*}$ ($t{=}14.7$) \\
MNIST & $n_h{=}256$, GN    & $88.0$ & $+0.71^{*}$ ($t{=}7.5$) \\
MNIST & $n_h{=}2000$, plain & $77.1$ & $+4.31^{*}$ ($t{=}12.5$) \\
MNIST & $n_h{=}2000$, GN    & $93.4$ & $+0.58^{*}$ ($t{=}10.1$) \\
\midrule
Fashion & $n_h{=}256$, plain & $56.8$ & $+1.39^{*}$ ($t{=}2.4$) \\
Fashion & $n_h{=}256$, GN    & $72.4$ & $-0.16$ ($t{=}-0.8$) \\
Fashion & $n_h{=}2000$, plain & $60.0$ & $+3.45^{*}$ ($t{=}9.6$) \\
Fashion & $n_h{=}2000$, GN    & $79.2$ & $+0.49^{*}$ ($t{=}4.5$) \\
\bottomrule
\end{tabular}
\end{table}

Two effects separate cleanly (Table~\ref{tab:twobytwo}).
(i)~\textbf{Normalization collapses the gain at matched width.}
Holding width fixed at $256$, inserting GN changes the $b{=}10$
learning-term gain from $+3.32$ to $+0.71$ on MNIST ($-2.61$) and from
$+1.39$ to $-0.16$ on Fashion-MNIST ($-1.55$); at width $2000$, from
$+4.31$ to $+0.58$ ($-3.73$) and from $+3.45$ to $+0.49$ ($-2.96$).
The effect is large, negative, and present at both widths.
(ii)~\textbf{Width moves the gain the other way.} Widening the plain
trunk $7.8\times$ \emph{increases} the gain ($+0.99$ on MNIST,
$+2.06$ on Fashion); under GN the width effect is small and mixed
($-0.13$, $+0.65$). Width therefore cannot explain the GN-MLP
collapse --- the wider plain network has, if anything, \emph{more}
repair gain, consistent with the headroom reading of
Sec.~\ref{sec:whattracks} (a wider trunk leaves more learning-term
room open). The decoupling conclusion survives the confound and is
stronger for it: the same-architecture control's verdict rests on the
normalization delta, not on the width delta it accidentally also
carried.

% =====================================================================


\section{What actually tracks the gain}
\label{sec:whattracks}

If not mortality, what? Table~\ref{tab:three} orders the six cells by
$\Delta\ell_{b=10}$: the bare MLPs lead on both datasets, the
strong-baseline configurations trail, and within MNIST the ordering
broadly tracks the unbiased baseline's learning-term weakness, though not
      monotonically
($\ell_0$ $74.3 \to 84.1 \to 93.5$ versus gain $+3.32 \to +0.07 \to
+0.57$; Fig.~\ref{fig:headroom}). We call this loose regularity
\emph{repair headroom}: the
bias can only buy learning-term room that the unbiased baseline
leaves open.

\begin{figure}[t]
\centering
\includegraphics[width=0.72\linewidth]{fig_headroom.pdf}
\caption{Repair gain at $b{=}10$ against the unbiased baseline's
learning term. The qualitative regularity (weak baseline, larger
gain) is visible within Fashion and is non-monotone within MNIST, where the
convolutional arm has the weakest baseline and the smallest gain, but a quantitative
law fails: ConvADNet-Fashion has $24.7$ points of $\ell$-headroom and
no significant gain. Headroom is necessary-looking, not sufficient.}
\label{fig:headroom}
\end{figure}

We stress what the data do \emph{not} support. A quantitative
headroom law fails: across the six configuration$\times$dataset
cells the Pearson correlation between HR$_\ell$ and $\Delta\ell$ is
$r = 0.38$ (cell means; pooled over per-seed differences, using each
cell's mean HR$_\ell$, $r = 0.27$; using each seed's own
$100-\ell_0$ gives $r = 0.31$), and the
 clearest counterexample is ConvADNet on Fashion --- $24.7$ points of
$\ell$-headroom, gain $+0.36$ ($t{=}1.4$, not significant). Headroom
is thus \textbf{necessary-looking but not sufficient}: a strong
baseline forecloses the repair (GN-MLP), but an open baseline does
not guarantee it (convolution adds an independent obstacle ---
plausibly the routing prior interacting with distributed conv
features, a question we leave open).

What is uniform across both strong-feature configurations is the
\emph{sign structure}: $\Delta\ell \ge 0$ (small) while
$\Delta\mathrm{avg} < 0$ (significant in $3/4$ cells;
Fig.~\ref{fig:tradeoff}).\hspace{0pt} The repair
never destroys the learning term, but on strong backbones it buys
that term by trading retention --- exactly the trade the bare MLP's identity
decomposition predicts, sharpened into a reversal once the learning
side has no room left to absorb the cost.

\begin{figure}[t]
\centering
\includegraphics[width=0.72\linewidth]{fig_tradeoff.pdf}
\caption{What the bias buys and what it charges
($T{=}10$, $b{=}10$, paired means). The bare MLP sits in the
gain--gain quadrant; every strong-feature cell falls in the
learning-vs-retention trade quadrant, with final accuracy
significantly negative in $3/4$.}
\label{fig:tradeoff}
\end{figure}

% =====================================================================



% ===== S4: 机制再定性 =====


% =====================================================================
\paragraph{The gate, in the notation used below.}
\label{par:gate}
The reference network is the one shared by the whole line of work
(ADNet): a one-hidden-layer active-dendrite network with $n_h{=}256$ hidden
units, one dendritic segment bank per task ($K{=}T$), $k$-winners-take-all
sparse output ($k{=}64$). Given input $x$ and context one-hot $c$:
\begin{equation}
\begin{split}
t &= W_{\mathrm{ff}} x, \qquad
u_k = w_k^{\top} c \;\; (k\text{-th segment row}), \\
h &= \mathrm{kwta}\!\left(t \odot \sigma(d)\right), \qquad
d = \max_k u_k,
\end{split}
\label{eq:adnet}
\end{equation}
i.e.\ the winning segment's scalar $d$ gates every hidden unit
multiplicatively. The repair studied in Sec.~\ref{sec:causal} adds a constant
$b$ to the own segment's $u_{\mathrm{own}}$ before the max, everywhere
(train, eval, usage).

\paragraph{Execution correction (2026-10-27).}
An independent data-consistency audit of this line, run after the first
version of the gate experiments had completed, found two defects in the
harness, and the experiments were re-run in full. (i)~The elastic penalty
was applied without its $\lambda$ factor, so every run had effectively
executed with $\lambda \approx 1$ while the record claimed
$\lambda{=}1000$; the penalty now carries $\lambda$ as the reference
implementation does. (ii)~The degenerate control had a hard-coded hidden
width of $2000$ while the dendritic arms use $256$, so the two arms of the
ablation differed by roughly eight-fold in parameter count; both now use
the same width. The numbers reported in Sec.~\ref{sec:mech} and
Sec.~\ref{sec:ablation} are all from the corrected re-run; the first
version's records are retained in the artifact under an explicit
deprecation notice. The correction moves the absolute accuracies
substantially --- a genuinely enforced penalty forgets more, and the base
network drops ${\sim}7.6$ points on MNIST --- and it changes one qualitative
verdict: the first version's MNIST gate-removal \emph{equivalence} was an
artifact of the width mismatch, and under matched widths the equivalence
becomes approximate rather than exact (Sec.~\ref{sec:ablation}). The
mechanism probes are unaffected.


\section{Preregistered design of the robustness probe}
\label{sec:design}

The design was locked on 2026-10-25, before any data was written to disk
(preregistration document included in the artifact). Noise is injected
\emph{at evaluation only}; training sees oracle contexts. The noise model
is worst-case in a specific sense: with probability $\rho$ per sample the
context one-hot is replaced by a uniformly random \emph{other} task ---
never the correct one --- so corrupted samples exercise the gate under a
plausible wrong context rather than a neutral one.

\begin{itemize}
\item \textbf{E1 (gating robustness).} $b{=}0$, $\rho \in
\{0, 0.1, 0.2, 0.4\}$, both datasets. Survival: $\rho{=}0.1$ costs
$\le 2$ points and $\rho{=}0.2$ costs $\le 5$, \emph{and} the
per-seed direction agrees across both datasets in $\ge 8/10$ seeds;
death: $\rho{=}0.1$ costs
$>5$ or $\rho{=}0.2$ collapses ($\bar a < 50$ or cost $>10$).
\item \textbf{E2 (repair survival).} Same grid crossed with
$b \in \{0, 10\}$, paired per $(\rho, \text{dataset}, \text{seed})$.
Survival: at $\rho{=}0.1$ the repair's paired $\Delta\bar a > 0$ and
retains $\ge 50\%$ of its $\rho{=}0$ gain; death: $\Delta\bar a \le 0$.
\item \textbf{E3 (soft mixture, exploratory).} Context replaced by
$(1-\rho)\,\mathrm{onehot} + \rho\,\mathrm{uniform}$; direction only,
no survival claim (multiple-comparison hygiene).
\end{itemize}

Two deviations from the pre-registered wording, both forced and both
reported rather than absorbed. First, E2's criterion was written for
$\Delta\ell$; the noise harness records only the final-phase
per-task accuracies and no $\ell$ field, so $\Delta\ell$ is not
computable from the collected data and the criterion is evaluated on
$\Delta\bar a$ instead --- the same "net benefit of the repair"
quantity, in the one form the runs support. Second, E1's third
condition (per-seed directional agreement) was not met in strict
form: the per-seed $\rho$-effect has mixed signs, though every
per-seed shift stays below $0.1$ points. We report that as a
condition we did not satisfy rather than dropping it, and the two
amplitude conditions pass with orders of magnitude to spare.

The noise stream is fixed per seed (a dedicated generator), so paired
comparisons share corruption patterns exactly. The protocol trains one
model per $(\text{dataset}, b, \text{seed})$ and evaluates it at all
four $\rho$ levels: \textbf{60 training runs} (40 for E1+E2, 20 for E3)
yielding 240 evaluation cells. (The pre-registration counted $\rho$ as
independent training runs --- 240 --- which over-counts by a factor of
four; the number of evaluated $(\rho, b, \text{dataset}, \text{seed})$
cells is unchanged.) A dated amendment documenting every deviation is
included in the artifact
(``prereg/S4\_ctxnoise\_预注册\_修订说明\_2026-10-27.md''); the original
pre-registration text is left untouched.



\section{Robustness to a wrong context: nothing degrades, at any noise level}
\label{sec:results}

\begin{table}[t]
\centering
\caption{Final average accuracy under worst-case context flips
($n{=}10$ seeds; noise applied to the final-phase evaluation only;
$\lambda{=}1000$ enforced, $n_h{=}256$ in all arms).}
\label{tab:main}
\begin{tabular}{llrrrr}
\toprule
dataset & arm & $\rho{=}0$ & $\rho{=}0.1$ & $\rho{=}0.2$ & $\rho{=}0.4$ \\
\midrule
MNIST    & base ($b{=}0$)    & 63.86 & 63.79 & 63.72 & 63.62 \\
MNIST    & repaired ($b{=}10$) & 68.67 & 68.67 & 68.67 & 68.67 \\
Fashion  & base ($b{=}0$)    & 42.20 & 42.15 & 42.09 & 42.04 \\
Fashion  & repaired ($b{=}10$) & 42.89 & 42.89 & 42.89 & 42.89 \\
\bottomrule
\end{tabular}
\end{table}

Table~\ref{tab:main} is the entire empirical content of E1 and E2.
Both preregistered survival criteria pass with orders of magnitude to
spare: the largest movement anywhere is $0.25$ points (MNIST base,
$\rho{=}0.4$), the largest on Fashion-MNIST is $0.16$, and the repaired
arm moves by at most $0.01$ points on MNIST and \emph{exactly} $0.000$
points on Fashion-MNIST --- where the death criteria allowed $10$ and
$5$. The repair's paired gain on MNIST is $+4.80$ ($t{=}17.0$) at
$\rho{=}0$ and $+5.05$ ($t{=}17.3$) at $\rho{=}0.4$; retention of the
gain is $100.0\% \to 105.1\%$ across the sweep --- under noise the
repair \emph{helps slightly more}, because its floor is flatter.
On Fashion-MNIST the gain is small, as the transfer test found at $T{=}10$
: $+0.69$ ($t{=}1.5$, n.s.) at $\rho{=}0$, rising
to $+0.86$ ($t{=}1.8$) at $\rho{=}0.4$ (retention $100\% \to 123.7\%$;
all n.s.). E1's third condition (per-seed directional agreement) holds
in strict form on MNIST --- $10/10$ seeds shift negatively at every
$\rho$ --- and fails on Fashion-MNIST ($8/10$, $9/10$, $7/10$ negative,
with every per-seed shift within $\pm 0.18$ points); we report it as
satisfied on one dataset and not on the other, as disclosed in
Sec.~\ref{sec:design}. E3's soft mixture tracks the oracle slightly
more closely than hard flips on the base arm (MNIST $\rho{=}0.4$: soft
$-0.09$ vs.\ hard $-0.25$; paired soft$-$hard $+0.16$, $t{=}3.8$) and
does nothing on the repaired arm; every E3 effect is ${\le} 0.26$
points.

A protocol-level artifact is worth excluding explicitly: the noise
\emph{is} reaching the network. Per-seed corruption patterns are fixed,
the context tensor is modified before the forward pass, and --- the
direct check --- swapping contexts \emph{does} change predictions in the
base network ($14.8\%$ of samples on average on MNIST, $30.1\%$ on
Fashion-MNIST, up to $51.7\%$ for particular swaps;
Sec.~\ref{sec:mech}). What is true is that prediction changes of this
kind average out across a $10$-way task mixture at the final-accuracy
level. The null is real, and it demands a mechanism.



\section{Mechanism: the gate does not gate by context}
\label{sec:mech}

Three probes, each applied to fully trained networks. The headline
numbers below are MNIST seed 0; the pattern replicates on Fashion-MNIST
seed 0 and MNIST seed 1 (all numbers in the artifact): at $b{=}0$,
flip rates $14.8$--$30.1\%$ mean / $23.8$--$51.7\%$ worst pair; at
$b{=}10$, flip rates $0.007$--$0.021\%$ mean / at most $0.15\%$ worst
pair; gate own/other identical to four decimals ($0.564$--$0.565$); and
mean column cosines in $[0.001, 0.003]$ in every case.

\paragraph{Prediction flip rate.} For each task $c$ and each alternative
context $c' \ne c$, measure the fraction of test samples whose argmax
prediction differs between $\mathrm{ctx}{=}c$ and $\mathrm{ctx}{=}c'$.
Base network: mean $14.8\%$ (MNIST) and $30.1\%$ (Fashion-MNIST), worst
pair up to $51.7\%$ --- the base output is genuinely, substantially
context-sensitive at the individual-prediction level, yet none of that
sensitivity survives to final accuracy (Sec.~\ref{sec:results}: at most
$0.25$ points). The context pathway does something per sample; what it
does nets out to nothing. Repaired network ($b{=}10$): mean
$0.007$--$0.021\%$, worst pair $0.15\%$ --- no longer \emph{exactly}
zero, but a rounding error: roughly one flipped prediction in five
thousand on average, a dozen samples for the worst single pair.

\paragraph{Gate response.} The winning segment's gate value
$\sigma(d)$ has mean $0.565$ under the correct context and $0.565$
under every alternative --- identical to four decimal places, in both
the base ($0.5651$) and repaired ($0.5652$) networks. The context
one-hot selects \emph{which} segment wins, but not \emph{how much} it
gates. As the control below shows, this value is not produced by
training either: it is already $0.564$ in the untrained network, so the
correct description is that the gate's value is a constant of the
\emph{initialization} of the gate bank, and a function of neither the
context nor the learning.

\paragraph{Column geometry.} The dendritic weights projected onto the
$T$ context columns remain mutually orthogonal after training: mean
pairwise cosine $0.003$ (per-model mean in $[0.001, 0.003]$; the most
negative single pair over the three models is $-0.056$). The columns are
in fact orthogonal \emph{before} training as well, and training barely
moves them: the cosine between each trained column and its own
initialization is $0.9995$ on MNIST and $0.9997$ on Fashion-MNIST
($n=3$ seeds). The orthogonality is therefore not something training
preserved; it is what the network started with.

\paragraph{Why the context pathway is unused: the gradient mask, not the
argmax.} We have been describing the constant gate value and the
orthogonal columns as properties of a trained network, and in the
previous version of this manuscript we attributed them to training. That
attribution was wrong, and the reason is worth stating precisely because
it is a defect of the architecture and its protection protocol rather
than of the learning problem.

The gate bank is a single matrix $W$ of shape
$(\text{nseg}\cdot\text{nh}) \times T$, laid out so that block
$(s, h)$ holds the row that segment $s$ uses for neuron $h$ and context
$t$. The forward pass computes $u = W\,\mathrm{ctx}$, and since
$\mathrm{ctx}$ is a one-hot, the value seen by task $t$ is
\emph{column} $t$ of $W$, in its entirety: all $\text{nseg}$ blocks of
it. The per-task gradient protection, however, is applied to the
\emph{blocks}: during the training of task $t$ only block $(t, \cdot)$ of
$W$ is allowed to receive a gradient. Block $(s, h)$ with $s \ne t$ is
therefore masked exactly when it is being read (it belongs to another
task's column) and also masked when its own task is trained (because
then block $(t,\cdot)$ is the one opened). The consequence is exact:
every off-diagonal block $W[s, \cdot, t]$ with $s \ne t$ receives a zero
gradient in every phase of the run and stays frozen at its random
initialization for the entire training.

We verified this directly rather than inferring it. Comparing each
trained network against a snapshot of its own initialization, the
relative displacement of the 90 off-diagonal blocks is
$0.0$ --- not small, but exactly zero --- in every seed on both
datasets, while the 10 diagonal blocks move by $6$--$10\%$; the trained
column agrees with its initialization to a cosine of $0.9995$. All
measurements in this section use $n=5$ seeds per dataset and were
reproduced on an independent machine to within $0.002$. The
weight matrix that the gate actually reads is, to five significant
figures, the weight matrix the network was born with.

Three consequences follow, and we separate them because they point in
different directions. First, the negative result of
Sec.~\ref{sec:anti} survives and is strengthened. If the absence of
task--segment specialization were an artefact of this mask, then
removing the mask should produce specialization. It does not: with the
mask switched off and the off-diagonal blocks free to move (they
displace by $7$--$11\%$), the own segment still wins only
$10.1\%$ of the neuron--sample slots on MNIST and $10.2\%$ on
Fashion-MNIST, against $10.1\%$ on both datasets in the untrained network
($n=5$ seeds per dataset, re-run on an independent machine).
The architecture does not learn task--segment routing under either
protocol, so the negative result is a property of the architecture.

Second, the constant gate value is not a training outcome. A freshly
initialized, untrained network already has a winning-segment gate value
of $0.5641$ on both datasets; after training it is $0.5646$
(MNIST) and $0.5644$ (Fashion-MNIST). These are $n=5$ seeds per dataset;
an independent re-run on a different machine reproduced every value in this
section to within $0.002$. The number the previous version of
this manuscript called "a constant of the training run" is a constant of
the \emph{initialization}, and it follows from the initialization
analytically: the winning row of a column is the maximum of
$\text{nseg}$ near-independent entries, so with the Kaiming-uniform
initialization of a bias-free $\text{fan}$-in-$T$ layer
($\sigma = 0.183$ at $T = 10$) the expected maximum is
$\mathbb{E}[\max_{s} z_s] \approx 0.28$, and
$\sigma(0.28) = 0.570$, which is the value we measure.

Third, and this is the uncomfortable one, the protection protocol is
not protecting what the method needs. The gradient mask was introduced
to keep each task's dendritic rows from drifting, and it does that
perfectly --- but because the gate reads whole columns, it also freezes
every entry of every column except one block, which is $90\%$ of the
context pathway. The mechanism this paper set out to audit is not merely
unused; it is, by construction, untrainable.

Together: the architecture's context pathway is present, protected per
task --- and, because the per-task gradient mask freezes $90\%$ of the
gate matrix, in practice untrainable as well as unused as a
\emph{modulator}. The max over segments is won by whichever row the
(orthogonal, near-random) columns happen to favor, the gate value is the
same constant ($\approx 0.565$, inherited from the initialization) no
matter which segment wins, and the repair adds a constant so large
($b{=}10$ against $|u| \lesssim 1$) that the own segment always wins
with $\sigma(10 + u) \approx 1$: a context-free gain of one. Equation~\ref{eq:adnet} degenerates, in both arms, to a
multiplication by a constant --- ${\approx}0.57$ in the base network,
${\approx}1.0$ in the repaired one. Which of those two constants
matters, and why, is what the next section measures.



\section{Ablation: what does the repair repair?}
\label{sec:ablstat}

\paragraph{Statistics.} Paired $t$-tests over $n=10$ seeds, Holm-corrected within each metric family across the three non-base arms. On MNIST the two substantive comparisons (repaired vs.\ base, gate-deleted vs.\ base) survive correction on both $\Delta\ell$ and $\Delta\mathrm{avg}$, with effect sizes $d_z=4.64$ and $2.46$ respectively --- the ordering of the two arms, not merely their significance, is what the mechanism section relies on. On Fashion-MNIST no comparison survives correction, which is consistent with the paper-wide finding that this dataset supports no effect of the repair at $T=10$.
\label{sec:ablation}

If the repaired network computes (approximately) $h = \mathrm{kwta}(t)$,
then a network with the dendritic bank deleted and the same trunk,
sparsifier, and anchor chain should reproduce most of the repair's
benefit --- and, conversely, a network whose gate is held at the base
system's measured constant ($0.565$) should reproduce the base system,
because the base gate already \emph{is} that constant
(Sec.~\ref{sec:mech}). We therefore run four arms under one protocol
($\rho{=}0$, $n{=}10$ paired seeds, all arms at $n_h{=}256$, all with
usage-weighted Fisher anchors at $\lambda{=}1000$): \textbf{base}
($b{=}0$), \textbf{repaired} ($b{=}10$), \textbf{no-gate}
($h = \mathrm{kwta}(W_{\mathrm{ff}} x)$; the bank deleted), and
\textbf{constant-gate} ($h = \mathrm{kwta}(0.565 \cdot
W_{\mathrm{ff}} x)$; the gate pinned at the base system's measured
mean). The base and repaired arms replicate the Sec.~\ref{sec:results}
runs to two decimals. The degenerate controls were added during the
post-audit re-run (Sec.~\ref{sec:setup}); they were not
preregistered. The no-gate control exists in three versions in the
artifact --- an invalidated first attempt (an extra sigmoid on $t$ that
changed the activation structure), a width-mismatched second version
($n_h{=}2000$), and the matched version reported here --- all retained
under an explicit deprecation notice (Table~\ref{tab:ablation}).

\begin{table}[t]
\centering
\caption{Four-arm control experiment ($n=10$ paired seeds, $\rho=0$, matched width $n_h=256$, matched usage-weighted anchors). Differences are paired against the base arm of the same dataset and seed; $p$ is Holm-corrected within each metric family; $d_z$ is the paired standardized effect size.}
\label{tab:ablation}
\resizebox{\textwidth}{!}{%
\begin{tabular}{llrrrrr}
\toprule
dataset & arm & $\Delta\ell$ & $p$ & $d_z$ & $\Delta\mathrm{avg}$ & $p$ \\
\midrule
fashion & const & $-0.33{}$ & $1.000$ & $-0.15$ & $+0.49{}$ & $0.969$ \\
fashion & repaired & $+1.39{}$ & $0.154$ & $+0.77$ & $+0.69{}$ & $0.526$ \\
fashion & nogate & $+1.02{}$ & $0.461$ & $+0.49$ & $+2.07{}$ & $0.146$ \\
mnist & const & $+0.07{}$ & $1.000$ & $+0.08$ & $-0.25{}$ & $0.969$ \\
mnist & repaired & $+3.32{*}$ & $0.000$ & $+4.64$ & $+4.80{*}$ & $0.000$ \\
mnist & nogate & $+2.61{*}$ & $0.000$ & $+2.46$ & $+3.64{*}$ & $0.000$ \\
\bottomrule
\end{tabular}
}
\end{table}


\textbf{The constant-gate control reproduces the base system, not the
repair.} Pinning the gate at $0.565$ lands at $63.62$ on MNIST (base:
$63.86$; paired difference $-0.25$, $t{=}-0.7$, n.s.) and $42.70$ on
Fashion (base: $42.20$; $+0.49$, $t{=}0.7$, n.s.). Holding the gate
constant buys nothing --- which is exactly what the mechanism section
predicts, since the base network's gate already \emph{is} constant at
that value. The entire difference between the base and repaired
networks is the value of that constant: $0.57$ damps the trunk,
$1.0$ does not.

\textbf{Gate deletion recovers most of the repair --- and on
Fashion-MNIST, more than the repair.} The no-gate control reaches
$67.50$ on MNIST, recovering $+3.64$ of the repair's $+4.80$-point gain
($76\%$) and falling short of the repaired network by $-1.16$
($t{=}-2.7$): keeping the (saturated) multiplicative structure is worth
about one point, marginally significant at $n{=}10$. On Fashion-MNIST
the no-gate control is the \emph{best} arm: $+2.07$ over base
($t{=}2.5$) and $+1.37$ over the repaired network ($t{=}1.9$, n.s.) ---
on the harder dataset, full deletion beats the repair outright.

\textbf{Sensitivity to anchor weighting.} A second no-gate run with
plain (unweighted) Fisher anchors --- a stronger, mismatched chain ---
lands at $63.63$ on MNIST and $41.10$ on Fashion, i.e.\ at or below
base. Anchor weighting alone moves the no-gate arm by $3.9$ and $3.2$
points. The matched-anchor four-arm table above is therefore the
primary comparison, and the residual MNIST gap ($-1.16$) should be read
against that sensitivity. The first version of this experiment reported
a clean MNIST \emph{equivalence} between the no-gate control and the
repair; that verdict was an artifact of the width mismatch disclosed in
Sec.~\ref{sec:setup} and is retracted.

\textbf{Verdict.} The repair is a saturating soft deletion: $b{=}10$
drives $\sigma \to 1$, which is the no-gate limit reached through the
gate rather than around it. What the repair buys on these benchmarks is
the removal of a constant ${\approx}0.57$ damping factor --- and on
Fashion-MNIST, where deletion is free to remove it \emph{exactly}, the
removal is worth slightly more than the repair itself. What does
\emph{not} follow is the stronger claim of the first version, that the
repair and gate removal are interchangeable on MNIST: with matched
widths and matched anchors they are close ($76\%$ of the gain
recovered) but distinguishable.




% ===== 自我诘问：dominant segment 是否为伪概念 =====


% =====================================================================
\section{Is ``dominant segment'' even a quantity? A self-interrogation}
\label{sec:pseudo}

The mechanism section (Sec.~\ref{sec:mech}) reports that the winning
segment's gate value is a constant of the initialization: $0.565$ under the
correct context and $0.565$ under every alternative, identical to four
decimals in both the base and the repaired networks, with the dendritic
columns remaining mutually orthogonal at their random initialization
(mean pairwise cosine $0.003$), and with $90\%$ of the gate matrix never
receiving a gradient at all. Taken literally, those measurements imply
something uncomfortable for the first half of this paper, and we would
rather state it than let a reviewer find it.

If the gate's value does not depend on which segment wins, then the
notion of a neuron's ``dominant segment'' is not a description of a
routing decision --- it is a description of an argmax over rows that are
nearly equal. And as we now know, those rows are nearly equal because
the protection protocol froze them: $90\%$ of the gate matrix never
receives a gradient, so the argmax ranges over a matrix that is still
essentially its random initialization. The near-degeneracy is not an
accident of the learning problem; it is built in. Under that reading, the quantity we spent
Sec.~\ref{sec:granularity} and Sec.~\ref{sec:audit} measuring --- the
probability that a neuron's dominant segment is its own task's segment ---
is a statistic of the \emph{instrumentation}, not of the network's
computation. The $0.05\to 0.01$ anti-specialization of
Sec.~\ref{sec:anti} would then not be a property of learned routing at
all, but a property of how argmax behaves on near-degenerate rows, and
the entire granularity ladder would be a null result about a null
quantity.

\subsection{What survives if that is right}

We checked this against the data rather than leaving it as a rhetorical
question, and the answer is that a substantial part of the paper survives,
while one specific claim does not. Three observations separate the
conclusions.

\textbf{The granularity ladder does not depend on the dominant-segment
instrument.} Its comparisons are \emph{within-architecture} paired
differences in the learning term $\ell$ and final accuracy between two
anchor constructions --- segment-resolved versus neuron-resolved --- under
identical training, seeds, and evaluation. No routing statistic enters
that comparison at all. The finding that segment-resolved anchors buy
nothing over neuron-resolved anchors ($|\Delta\ell|\le 0.24$ across every
dataset, depth, and elastic weight) is a statement about where the
\emph{benefit of usage localization} is realized, and it would survive
even if the dominant-segment measure were discarded entirely.

\textbf{The dead-neuron artifact does not depend on it either.} The
artifact (Sec.~\ref{sec:artifact}) is a statement about what happens to
an argmax instrument when $56$--$63\%$ of the rows it operates over are
exactly zero: the tie-breaking rule assigns them all to segment $0$.
That is a fact about argmax on a matrix with a zero block, verifiable
without any claim about what the network learned, and its remedy
(active-only filtering) is likewise an instrumentation change. If
``dominant segment'' is a pseudo-quantity, this finding becomes
\emph{more} important rather than less: it becomes a warning that a
broad class of measurements in this literature --- not only ours --- may be
reading a tie-breaking artifact.

\textbf{The mechanism and the ablation are about the gain, not about
routing.} Sections~\ref{sec:mech}--\ref{sec:ablation} measure the
context-flip rate, the gate value, the column geometry, and the final
accuracy of four arms. Every one of those is a measurement of the
network's output or of an intervention's effect on that output. None of
them is a measurement of ``specialization''.

\subsection{What does not survive}

One claim does not survive, and it is the one we would have made most
strongly. We cannot assert that segment--task specialization is
\emph{actively violated} as a property of the learned computation, because
if the gate does not select by context, there is no learned
segment--task mapping to violate. The honest formulation is weaker and
we adopt it throughout: \emph{no evidence of segment--task specialization
is detectable by the standard instrument, and the instrument itself is
unreliable in this architecture family}. The distinction is between a
measured absence and an absence of measurement, and after the mechanism
section we can only claim the former.

\subsection{Why this is the useful reading}

The weaker formulation is not a retreat; it is the correct one, and it
sharpens what the paper is actually for. Three things survive in full
and are not affected by the pseudo-quantity worry.

(i)~The granularity question --- \emph{can protection usefully descend
below the neuron?} --- is answered by comparing anchor constructions, and
the answer is no, for a reason we pin mechanistically (attribution is
self-limiting; segment-level ownership lacks a separable carrier). This
is the paper's primary result and it never required a routing statistic.

(ii)~The instrument warning stands on its own. A measurement that
fabricates $60\%$ dominance from a zero block is a hazard to any study
that uses it, and the hazard is demonstrated by construction rather than
by comparison to a ground truth we do not have.

(iii)~The mechanism result --- that the own-segment bias works by
saturating a gate that was never context-dependent, and that deleting the
gate reproduces the gain --- is a statement about interventions and
outputs. It is also the most transferable finding in the paper, because
it does not presuppose that the architecture routes at all.

We note in closing that this self-interrogation was not in the original
version of this line of work, and that the instrument it produced --- a
gate-value probe with paired context flips --- was cheap. The reason it
came last rather than first is that the same thing that made the
apparent repair look convincing --- a large, consistent, saturating
dose--response --- also made it seem unnecessary to ask what the gate was
doing. That is a failure mode worth naming: \emph{a result strong enough
to be interesting is strong enough to be worth explaining}, and we had
the instrument to explain it before we thought to build it.



% ===== 相关工作 / 局限 / 讨论 =====


\section{Controls added after the first review of this line of work}
\label{sec:controls}

The five sections above are the argument as it stood when this manuscript
was assembled. A review of that draft pointed out five places where a
control was missing rather than a claim being wrong, and we ran them. Four
of the five changed what we can claim; one of them changed the paper's
relationship to the method it audits. We report them here, in the order
they were asked for: the specificity of the intervention
(Sec.~\ref{sec:ctrl-specificity}), a positive control for the instrument
(Sec.~\ref{sec:ctrl-oracle}), the construction-diversity of the
granularity result (Sec.~\ref{sec:ctrl-anchor}), the comparison with
contemporary protection methods (Sec.~\ref{sec:ctrl-baselines}), whether
the robustness null is a metric artefact (Sec.~\ref{sec:ctrl-metric}), and
the equivalence testing that replaces our earlier reliance on failure to
reject (Sec.~\ref{sec:ctrl-tost}). A closing paragraph
(Sec.~\ref{sec:ctrl-unchanged}) states what all of this leaves intact. We
have edited the earlier sections and the abstract to reflect the outcomes
rather than leaving the corrections implicit.

\subsection{The own-segment specificity control, and what the bias is
actually doing}
\label{sec:ctrl-specificity}

If the own-segment bias works by restoring own-segment routing, then the
word \emph{own} is doing work: the same bias applied to a different
segment should do less. If instead it works by removing a constant
damping factor, the qualifier should be irrelevant. We ran four arms on
the identical protocol (Permuted-MNIST and Permuted-Fashion, $T{=}10$,
$n_h{=}256$, $\lambda=1000$, usage-weighted Fisher anchors, $n=10$ paired
seeds): \emph{base} (no bias), \emph{repaired} (the bias of $b=10$ on the
own segment, as in Sec.~\ref{sec:causal}), \emph{random} (the same $b=10$
on one segment drawn at random per neuron and held fixed for the run),
\emph{uniform} (the same $b=10$ on every segment), and \emph{rescale}
(no gate at all: $h=\mathrm{kwta}(t/0.565)$, a pure constant removal of
the damping factor); Table~\ref{tab:ctrl-specificity} reports all five.

\begin{table}[htbp]
\centering\small
\caption{Repair gain against neuron mortality in the width$\times$normalization
factorial scan ($n=10$ paired seeds per cell). The gain is the paired
difference in final accuracy between $b{=}10$ and $b{=}0$; mortality is the
dead-neuron fraction of the repaired arm at the end of the sequence, measured
on the raw, bias-free network. Across the eight cells the two are
uncorrelated (Pearson $r=-0.11$, $p=0.79$; Spearman $\rho=+0.05$,
$p=0.91$), and the ordering inverts: the cells with the lowest mortality have
the smallest repair gain.}
\label{tab:scanmort}
\resizebox{\textwidth}{!}{%
\begin{tabular}{llrrrrrrr}
\toprule
dataset & width & GroupNorm & $n$ & mortality ($b{=}0$) & mortality ($b{=}10$) &
$\Delta$ final accuracy & $t$ & $p$ \\
\midrule
mnist & 256 & no & 10 & 0.584 & 0.535 & $+4.67$ & $+15.7$ & $<10^{-6}$ \\
mnist & 256 & yes & 10 & 0.004 & 0.000 & $+0.95$ & $+6.2$ & $<10^{-3}$ \\
mnist & 2000 & no & 10 & 0.945 & 0.939 & $+4.62$ & $+12.5$ & $<10^{-4}$ \\
mnist & 2000 & yes & 10 & 0.745 & 0.780 & $-2.74$ & $-2.3$ & $0.044$ \\
fashion & 256 & no & 10 & 0.618 & 0.592 & $+0.90$ & $+2.1$ & $0.063$ \\
fashion & 256 & yes & 10 & 0.010 & 0.020 & $+0.42$ & $+1.1$ & $0.321$ \\
fashion & 2000 & no & 10 & 0.952 & 0.948 & $+2.13$ & $+3.4$ & $0.008$ \\
fashion & 2000 & yes & 10 & 0.875 & 0.884 & $-6.64$ & $-4.0$ & $0.003$ \\
\bottomrule
\end{tabular}%
}
\end{table}
\FloatBarrier

\begin{table}[t]
\centering\small
\caption{Where the own-segment bias acts. Means over $n=10$ paired seeds on the identical protocol. The ``random arm places the same bias on a segment drawn at random per neuron; the ``uniform arm places it on every segment; the ``rescale arm removes the gate entirely and divides the trunk output by the base network\'s measured gate constant $0.565$.}
\label{tab:ctrl-specificity}
\begin{tabular}{llrrrr}
\toprule
dataset & arm & $\ell$ & final avg & BWT & $n$ \\
\midrule
mnist & base & 74.38 & 63.83 & -11.73 & 10 \\
mnist & repaired & 77.52 & 68.55 & -9.97 & 10 \\
mnist & random & 74.19 & 64.04 & -11.29 & 10 \\
mnist & uniform & 77.38 & 68.28 & -10.11 & 10 \\
mnist & rescale & 81.23 & 73.63 & -8.44 & 10 \\
fashion & base & 56.83 & 42.40 & -16.02 & 10 \\
fashion & repaired & 57.76 & 43.16 & -16.22 & 10 \\
fashion & random & 56.62 & 42.39 & -15.82 & 10 \\
fashion & uniform & 58.23 & 43.66 & -16.19 & 10 \\
fashion & rescale & 61.50 & 47.84 & -15.19 & 10 \\
\bottomrule
\end{tabular}
\end{table}

The result sharpens the story in a way we did not expect. On MNIST the
\emph{repaired} arm reaches $\ell=77.5$ against the base arm's $74.4$,
and the \emph{uniform} arm --- the same bias on every segment rather
than the own one --- reaches essentially the same place ($77.4$). So
\emph{own} is not doing the work: a bias on the own segment is worth
about three points of learning term, and a bias on \emph{any} segment is
worth the same. What does not work is the \emph{random} arm, which
reaches the base arm's level ($74.2$, paired difference $-0.19$,
$p=0.34$): biasing a segment that the network does not use buys
nothing, and biasing a segment it does use buys the full effect,
whichever one that is. The mechanism is therefore about which segment
is active, not about the segment being the task's own.

Two consequences follow. The first is about the paper's own vocabulary:
we have been describing the intervention as an own-segment bias
throughout, and the controls say the accurate description is a bias on
\emph{whatever segment the network is currently using}. The second is
about the size of the effect: the \emph{rescale} arm, which removes the
gate altogether and divides the trunk output by the base network's
measured gate constant $0.565$, reaches $\ell=81.2$ --- four points
above the bias arms. We wrote in Sec.~\ref{sec:mech} that the bias buys
the removal of a constant damping factor; this control says the bias
buys that, and that a more direct way to buy it does better still. The
own-segment bias is best understood as a crude, one-parameter way of
removing a damping factor that the architecture should never have
introduced, not as a repair of routing.

\subsection{A positive control for the instrument we say is unreliable}
\label{sec:ctrl-oracle}

Sec.~\ref{sec:audit} rests on a negative reading of the dominant-segment
instrument, and Sec.~\ref{sec:pseudo} argues that this reading may be a
statement about the instrument rather than about the network. A
measurement that only ever returns ``not detected'' cannot distinguish
``the quantity is absent'' from ``the instrument is broken''. We therefore
built a model in which task-dominant gating is true by construction: the
own segment's gate entry is forced to $+M$ and every other segment's to
$-M$, so that the own segment wins every competition by construction
($M=5$), with a weaker version ($M=1$) to locate the detection
threshold, and the base network unchanged. The override is applied at
the gate itself, so that both the network's forward pass and the
instrument read the forced values.

The instrument detects the forced specialization exactly: the
active-only own-segment rate is $1.000$ in every phase of every oracle
run on both datasets, against $0.055$ (MNIST) and $0.057$ (Fashion) for
the base network. The instrument is therefore sensitive, and the negative
readings in Sec.~\ref{sec:audit} are informative: in this architecture
family, task-dominant gating is genuinely not there. This is the
strongest support we can offer for the paper's central negative claim,
and it is support we did not have in the first draft.

The oracle arms are also informative about a different question, though
the answer is more mixed than we expected. Forcing the specialization
does improve the learning term on MNIST ($\ell=77.3$ against the base
network's $74.4$, paired $t(9)=11.5$, $p<10^{-4}$ after correction) and
slightly on Fashion-MNIST ($\ell=58.0$ against $56.8$, $t(9)=2.5$,
$p=0.034$ before correction and not significant after it). So a
specialization \emph{would} be useful on MNIST if the architecture could
acquire one. That is a statement about the architecture's capacity rather
than about the protection method, and it lines up with what
Sec.~\ref{sec:ctrl-specificity} finds: what the network needs is not
more protection but a usable path for the signal it already computes.

\subsection{The granularity result is a property of our anchor
construction, not of granularity as such}
\label{sec:ctrl-anchor}

The null in Sec.~\ref{sec:granularity} --- segment-resolved anchors buy
nothing over neuron-resolved ones --- was attributed to two properties of
the construction we chose: attribution is \emph{self-limiting} (the
anchor's own penalty pins the weights, so the tracked reference never
moves), and the Fisher mass that drives the weighting is too diffuse to
partition. A reviewer can reasonably object that the first of these is a
property of one implementation rather than a general fact about
protection below the neuron level, and that the null is therefore close
to a strawman. We tested it by building a second segment-resolved
construction that is not self-limiting: the per-segment reference
states are refreshed at the end of each phase in which that segment is
trained, so the penalty no longer pulls the weights toward a stale
target (variant V2; a third variant, V3, additionally estimates the
Fisher term per segment from the samples that segment wins).

The null does not survive. On MNIST, the original construction reproduces
it exactly ($\Delta\ell=+0.05$ on MNIST and $-0.04$ on Fashion-MNIST
against the neuron-resolved anchor, $p=0.76$ and $p=0.87$), while the
refreshed-reference variant reaches
$\Delta\ell=+2.49$ on MNIST ($t(9)=14.3$, $p<10^{-5}$ after correction,
$d_z=4.5$) and $+1.82$ on Fashion-MNIST ($t(9)=4.6$, $p=0.0013$), and the
segment-selective-Fisher variant reproduces it exactly. The general claim that sub-neuron protection
granularity is worthless is therefore \emph{not supported by our data},
and we withdraw it.

What the variant buys is not free, and the second number matters as much
as the first: the same variant that gains two and a half points of
learning term on MNIST and nearly two on Fashion-MNIST loses about twenty
points of final accuracy on MNIST ($\Delta\mathrm{avg}=-19.7$,
$t(9)=-25.8$) and nine on Fashion-MNIST ($-9.5$, $t(9)=-17.4$). Refreshing the per-segment
references lets a segment\'s anchor follow the segment rather than
references lets a segment's anchor follow the segment rather than
pinning it, which is exactly what makes the segment-level penalty
informative --- and exactly what makes it interfere with the tasks that
follow. So the honest summary of this control is not that
sub-neuron granularity is valuable, but that \emph{it is a real lever
whose direction depends entirely on the reference-update rule, and that
the null we reported was an artefact of the rule we happened to
implement}. A method that can be made to help or to hurt by a
implementation choice is not a settled question, and the first draft of
this paper treated it as one.

\subsection{How the method we audit compares with contemporary
alternatives}
\label{sec:ctrl-baselines}

Every comparison in this manuscript so far has been against elastic
regularization from 2016--2018 (EWC, SI, MAS) and PackNet. That is a
real weakness: the reader has no way to know whether the protection
scheme we are auditing is competitive with what the field does now. We
ran two contemporary replay-based methods on the identical protocol
($T=10$, $n_h{=}256$, $n=10$ paired seeds, same permutations, same
evaluation): \emph{ER}, a reservoir buffer of $500$ exemplars mixed into
each batch, and \emph{DER}, the same buffer augmented with the stored
logits and a distillation term. Neither uses the elastic anchor chain,
because replay is an alternative protection mechanism rather than an
addition to one --- stacking the $\lambda=1000$ penalty on top of a
buffer would suppress the gradient the buffer exists to supply.

The comparison is not close. On Permuted-MNIST, DER reaches a final
average accuracy of $86.4$ against DeLP's $64.0$, with a forgetting
measure of $-2.9$ against DeLP's $-11.4$; on Permuted-Fashion, $75.3$
against $42.5$, with $-0.7$ against $-15.9$. Every one of the twelve
paired comparisons is significant at $p<10^{-6}$ with standardized effect
sizes between $2.2$ and $14.0$, in DER's favour on all three metrics.
Experience Replay alone is also ahead of DeLP on final accuracy, though
it forgets more on MNIST than DeLP does.

This does not contradict anything in the previous sections, and it is
worth being precise about why. Every finding above is a statement about
the internals of DeLP and of the active-dendrite architecture it is built
on: that the protection granularity below the neuron buys nothing under
our anchor construction, that the routing statistic the field measures
is contaminated, that the own-segment bias works by removing a damping
factor, and that the whole scheme is outcompeted by a replay buffer
coupled with logit distillation. Those findings say that the particular
protection mechanism examined here is not competitive, and that the
architectural premise it inherits --- that tasks occupy different
dendritic segments --- does not hold in the trained network. They do not
say that the architecture is unprotectable. A reviewer should read this
paper as an argument that the dendritic-segment account of continual
learning, and the protection scheme built on it, need to be rethought
from the trunk upward, not as an argument that continual learning on this
backbone is hard.

\subsection{Is the robustness null a real null or a metric that cannot
see?}
\label{sec:ctrl-metric}

Sec.~\ref{sec:results} reports that nothing degrades under context
noise, while also reporting that the base network's predictions flip
under a context change at rates of $14.8\%$ (MNIST) and $30.1\%$ (Fashion)
at the individual-sample level. A reader is entitled to ask how both can
be true, and whether the accuracy metric is simply too coarse to see a
damage that is happening. We quantified it. The accuracy loss is linear
in the corruption rate $\rho$ with a slope of $-0.77$ points per unit
$\rho$ on MNIST (base arm), so the largest corruption in the sweep costs
$0.25$ points, and the $95\%$ confidence interval on that loss excludes
zero. Against the observed argmax flip rate of $14.8\%$, that slope
implies that only about $4$--$5\%$ of the flips turn a correct prediction
into an incorrect one --- the rest turn an incorrect prediction into a
correct one, or into another incorrect one, and these two effects
largely cancel under the 10-way task mixture. To lose five accuracy
points at $\rho=0.4$ the slope would have to be $12.5$ points per unit
$\rho$; the measured slope is $0.77$, a factor of sixteen smaller. The
null is therefore a real null: the noise reaches the network, the metric
is sensitive enough to see a damage of the size the pre-registered
thresholds were written for, and the damage does not occur.

\subsection{Equivalence testing for the comparisons we describe as
showing no difference}
\label{sec:ctrl-tost}

Several comparisons in this manuscript are described as showing no
difference: the granularity ladder, the constant-gate control against the
base network, and the repaired arm against the base arm on
Fashion-MNIST. A failure to reject is not evidence of equivalence, and at
$n=10$ a 95\% confidence interval on the learning term is wide enough
that a half-point difference cannot be excluded. We therefore re-ran each
of these as a two-one-sided-tests equivalence test at a pre-declared
margin of equivalence of $0.5$ accuracy points, reported as the 90\%
interval that corresponds to $\alpha=0.05$ in a TOST. The full table is
in the supplementary material.

The outcome is mixed and we report it as measured. For the granularity
ladder, the branch-versus-neuron comparison on the learning term is
equivalent at the declared margin in both MNIST cells, but not in the
Fashion cells and not for final accuracy; the correct summary is that we
detect no difference larger than roughly half a point where the effect
is small, not that the two constructions are interchangeable everywhere.
For the constant-gate control, no comparison meets the equivalence
criterion: on Fashion-MNIST the 90\% interval on the learning term is
$[-1.58,+0.92]$, so an effect of up to about one and a half points cannot
be excluded. We therefore describe that arm in the text as being
consistent with the base system rather than as reproducing it, and the
same applies to the repaired arm on Fashion-MNIST.

\subsection{What the additional controls leave unchanged}
\label{sec:ctrl-unchanged}

Two things survive all of this. The dead-neuron artifact is a fact about
argmax over a block of exactly-zero rows, verifiable without any claim
about what the network learned, and the corrected instrumentation it
motivates is an instrumentation change. And the self-interrogation of
Sec.~\ref{sec:pseudo} is unaffected: if the gate's value does not depend
on which segment wins, the dominant-segment statistic is a statistic of
the instrumentation, and the oracle control in Sec.~\ref{sec:ctrl-oracle}
shows the instrument is sensitive --- which means the negative readings
are real measurements of an absence, not artefacts of a broken measure.
What the additional controls remove is the strong form of the paper's
first contribution, the specificity of the own-segment framing, and any
implication that DeLP is competitive with contemporary protection
methods.


\section{Related work}
\label{sec:related}

\paragraph{Active-dendrite networks.} Single-neuron compartmental models
\citep{poirazi2003,london2005,poirazi2020,ujfalussy2018} motivate
multi-segment artificial neurons; Iyer et al.\ \citep{iyer2022} operationalize
context-gated segment competition for lifelong classification and report
reduced interference, and the GWC route-suppression variant
\citep{boyle2025,baek2024} and neuromorphic implementations
\citep{boahen2022} extend the family. These works validate the architecture;
to our knowledge none measures whether the intended segment--task
specialization actually obtains in the trained model, which is the object
of our audit --- nor, consequently, the dead-neuron artifact that
naive dominance statistics incur under production sampling budgets.

\paragraph{Elastic weight consolidation and its localizations.} EWC
\citep{kirkpatrick2017} protects parameters with a single global anchor;
SI \citep{zenke2017}, MAS \citep{aljundi2018}, and online variants
\citep{chaudhry2018} refine the importance estimate; masking approaches
\citep{mallya2018,serra2018,masse2018} hard-allocate capacity. \delp{}
\citep{zhou2026delp} localizes EWC per neuron via usage. A companion
study \citep{zhou2026anchorbudget} diagnoses the anchor-accumulation
problem --- the protection penalty grows with sequence length --- and
bounds the anchor stack. The present study is orthogonal to that one, and
the division of questions is worth stating explicitly, because a reader
who takes both as a single program could reasonably expect them to agree
rather than to answer different things. The companion study asks
\emph{how large a protection set should be}: it treats the anchor stack
as the object, bounds the accumulated Fisher mass, and derives a
log-linear law relating the learning term to the protected-mass horizon
($\ell \approx 84 - 11\ln h$, per-seed $R^2 \ge 0.98$, replicated at a
second task count and on Split-CIFAR-100 with a ResNet-18 backbone). This
study asks a different question: \emph{what should a protection set
protect?} --- in particular, whether protection granularity can usefully
descend \emph{below} the neuron. It answers negatively, and then explains
the negative via a routing audit. The two results are compatible and
neither implies the other: bounding the stack tells you how much
protection to keep, not whether finer protection is worth having, and a
bounded stack of neuron-resolved anchors remains the right object even if
segment-resolved anchors would have been better.

\paragraph{Negative results and measurement validity.} Benchmark pitfalls
\citep{musgrave2020,lucic2018,melis2018,agarwal2021,lipton2019} established
that evaluation practice can manufacture effects; the dead-neuron artifact
here is an instance inside the model rather than the benchmark: the
instrumentation, not the metric, fabricates the signal.

% =====================================================================


\section{Discussion: boundary conditions of the repair}
\label{sec:boundary}

\paragraph{For practitioners.} An own-segment bias is cheap,
architecture-agnostic, and --- on weak baselines with hard routing ---
effective . This paper draws its boundary:
\begin{itemize}
\item \textbf{Audit mortality before trusting any dominant-segment
statistic.} The bare-MLP artifact appears iff gates die; conv+GN mortality
is zero and the statistic reads clean. One probe (an untrained-pass
share, or the dead-gate fraction) is enough to know which regime you
are in.
\item \textbf{Audit baseline strength before deploying the repair.}
If the unbiased baseline's learning term is already near its ceiling,
the bias will not buy learning --- but will still charge retention
($\Delta\mathrm{avg}<0$ on all four strong-feature cells (from $-0.84$ to $-6.55$),
significant in $3/4$).
\item \textbf{Protection strength does not transfer.} $\lambda$
calibrated on one parameter scale left the conv baseline with
$2$--$4.4\times$ the MLP's forgetting. Any cross-architecture
comparison of protection methods must re-calibrate or report the
uncalibrated gap explicitly.
\end{itemize}

\paragraph{For the bare-MLP story.} The companion paper's chain survives
intact on its home architecture, and this paper sharpens its
mechanism claim: the repair was never a routing \emph{prior} in the
strong sense (the bare MLP's task-relevance probe already showed
zero-relevance tasks are untouched at every dose); it is a structured capacity
intervention whose raw material --- weak-baseline, dying-gate
networks --- is architecture-dependent. ``Structured capacity with an
architecture-bound supply.''

\paragraph{Threats to validity.} (i) The conv $\lambda$ was not
re-calibrated; a per-architecture sweep might restore some repair
effect, though the GN-MLP control --- where the MLP-calibrated
$\lambda$ \emph{is} correct --- shows the same collapse, so
mis-calibration is not the whole story. (ii) One conv topology; the
zero-mortality finding is a property of GroupNorm's affine output and
should replicate across conv trunks, but transformer backbones with
routing modules remain untested. (iii) Six configuration$\times$dataset
cells bound how strongly any cross-configuration regularity can be
asserted; we deliberately report the failed quantitative version
(Sec.~\ref{sec:whattracks}) rather than promoting the qualitative one
beyond its evidence.

% =====================================================================


\section{Discussion: what a repair of this shape is actually doing}
\label{sec:repair-disc}

\paragraph{For the bare-MLP story.} The dose--response audit established that
the own-segment bias repairs the learning term; the signed-drift measurement showed the network never pushes back at the parameter level.
The mechanism section supplies the missing re-qualification: the repair does not
restore routing --- there was no routing to restore --- it \emph{saturates
the gate into a context-free gain of one}, i.e.\ a soft deletion of a
constant ${\approx}0.57$ damping factor. Read together, the chain is
now: no specialization (the bare-MLP arm) $\to$ repair works (the bare-MLP arm) $\to$ the repair is
a saturating gate deletion that recovers $76\%$ of the removal benefit
on MNIST and is beaten by full removal on Fashion-MNIST $\to$ the repair does not transfer to strong backbones. What
survives as the ``mechanism'' of the repair is the removal of a constant
damping --- sparse capacity management at best, context-gated routing
not at all.

\paragraph{For evaluation practice.} Robustness certificates obtained
under context corruption are usually read as good news. In this family
they are a diagnostic: a system whose accuracy is invariant to
\emph{adversarial} context corruption is a system whose context pathway
is not load-bearing. The three-instrument protocol
(Sec.~\ref{sec:results}--\ref{sec:ablation}) costs a few GPU-hours and
converts that ambiguity into a verdict. We suggest it as a standing
companion to any context-gated continual-learning evaluation, in the
same spirit as the dead-neuron probe of the bare-MLP arm.

\paragraph{Limitations.} (i) Noise is injected at evaluation only;
training-time context errors corrupt usage statistics and anchors
differently, and are the natural next experiment. (ii) The mechanism
probes are computed on fully trained networks at $T{=}10$; whether the
context columns stay orthogonal at larger $T$ or longer training is
untested. (iii) The degenerate controls were not preregistered and were
added after the audit; the anchor-weighting sensitivity of the no-gate
arm ($\pm 3.9$ points, Sec.~\ref{sec:ablation}) is larger than the
residual repair-vs-deletion gap it bounds. (iv) One architecture
family, two digit datasets --- the same scope limits as the companion
line.



\section{Limitations}
\label{sec:limitations}

\begin{enumerate}
\item \textbf{Scale.} All results are on permuted benchmarks with $T \le 20$; the trunk is a
      one-hidden-layer MLP except in the convolutional and width-variant arms
      of Secs.~\ref{sec:conv}--\ref{sec:gnmlp}. The artifact and
      anti-specialization findings concern the interaction of kWTA
      mortality with argmax instrumentation and should be re-checked
      wherever segments exist (e.g.\ convolutional or transformer backbones
      with routing modules); the causal repair's transfer is untested
      beyond this family. The task-relevance probe
      (Sec.~\ref{sec:taskrel}) shows the repair is a structured capacity
      effect rather than a routing prior: it lifts in-domain tasks above
      chance and compresses task-to-task variance, but adds nothing for
      zero-relevance (random-label) tasks --- so its benefit is
      contingent on the task stream containing learnable structure, a
      condition that should be verified per architecture and benchmark
      before generalizing.
\item \textbf{The repair is a prior, and the prior bypasses the gate.}
      The constant bias
      encodes the correct assumption (own segment first) that the network
      fails to learn at effective strength. We measured this directly:
      given a trainable per-(task, segment) bias table initialized at
      zero, the network acquires only a small positive own-segment bias
      ($+0.113 \pm 0.031$ MNIST, $+0.083 \pm 0.025$ Fashion; $n{=}7$
      runs each, $t(6) \ge 8.7$) --- the raw-gate scale, an order of
      magnitude below the smallest effective dose --- and its accuracy
      matches the unbiased baseline. The mechanism section (Secs.~\ref{sec:mech}--\ref{sec:ablation}) establishes what the bias actually
does: it saturates the own-segment sigmoid until the gate stops
      depending on the context (context-swap flip rate $0.007$--$0.021\%$ at
      $b{=}10$), and on MNIST a gate-removal control recovers $76\%$ of the
      gain. The repair therefore does not restore segment--task routing;
      it removes the network's dependence on it. On Fashion-MNIST the
      comparison is not significant at $n{=}10$ ($t{=}1.9$, $p_{\mathrm{Holm}}=0.15$),
      so the mechanism is dataset-dependent. The routing-restoration reading
      is withdrawn outright by Secs.~\ref{sec:mech}--\ref{sec:ablation}:
      what the bias buys is the removal of a constant damping factor,
      and on MNIST a gate-removal control reproduces the gain outright.
      We also tested scheduled variants
      (Suppl.~S4): an annealed bias ($b_p = 10(1-p/T)$) retains
      $91$--$103\%$ of the constant-bias learning gain at $T{=}20$ but
      does not reliably remove the BWT cost (significant improvement in
      one of two datasets only), while a ramped bias ($b_p = 10p/(T{-}1)$)
      \emph{exceeds} the constant gain everywhere ($123$--$197\%$
      retention, $t \ge 2.6$) at an equally or more negative BWT cost.
      Scheduling moves along the retention--plasticity frontier ---
      tracking cumulative bias exposure, not timing --- but does not
      move the frontier itself; separating gain from cost appears to
      require a mechanism (e.g.\ task-staggered protection), not a
      schedule.
\item \textbf{Depth-dependent forgetting trade-off.} At $T{=}20$ the BWT
      cost of the bias is small but real: $3/10$ cells significantly
      negative ($-0.66$ to $-0.95$), the rest indistinguishable from
      baseline. We
      report it rather than select it away; the identity
      Eq.~\eqref{eq:identity} makes explicit that accuracy gains can
      coexist with mildly worse retention.
\item \textbf{Granularity null is architecture-bound.} The
      plain$\to$neuron step's dominance may weaken in architectures where
      segment parameters are genuinely disjoint (no shared feed-forward
      weights); our null is about this family's shared-substrate design.
\end{enumerate}

% =====================================================================


% =====================================================================
\section{Conclusion}
\label{sec:conclusion}

We asked whether the segment structure of active-dendrite networks offers
protection granularity below the neuron, and found that the question
dissolves one level up: the entire benefit of usage-localized protection is
realized at neuron level because the segment--task specialization the finer
granularity would need is not detectable by the standard instrument --- at
chance before each phase and pushed below chance by that phase's own training,
with Fisher mass too diffuse to partition meaningfully. Because that
instrument is unreliable in this family, this is an absence of measurement
rather than a measured absence, and we claim no more than that. On the way we isolated a dead-neuron artifact that
fabricates dominance statistics for any naive instrument in this family.

Having measured the size of the gap, we tried to close it. A one-parameter
bias on the own-segment dendritic input does raise the learning term in every
tested cell (significant in
$16/22$ dose cells, $\Delta\ell$ up to $+4.6$,
$t(9){=}8.1$), with final accuracy improving in step. We had reported this as
a causal repair, and it is not one --- not because the measurement is wrong,
but because the interpretation was. Three further results say so.

The repair does not transfer. On a convolutional active-dendrite backbone the
dose--response collapses to noise and final accuracy reverses; the natural
explanation (the repair revived dying gates, and conv gates do not die) is
refuted by a same-architecture control in which mortality and gain do not track
in both directions. What the repair actually needs is a weak baseline with
learning-term room to buy; what it always charges is retention, and on
strong-feature backbones --- where nothing else is for sale --- that is all it
buys.

And the mechanism is not the one we assumed. The context pathway in these
networks is present, protected per task, and unused as a modulator: the
winning segment's gate value is a constant of the initialization ($0.564$
before any training), identical under every alternative context, and the
dendritic columns remain orthogonal at their random initialization. We
further identify why: the per-task gradient mask freezes every
off-diagonal block of the gate matrix at exactly its initial value, and
because the forward pass reads whole columns rather than blocks, $90\%$
of the context pathway never receives a gradient in the entire run. The
pathway is therefore untrainable, not merely untrained --- although
removing the mask does not make specialization appear, which is why the
negative result stands either way. At $b{=}10$ the own segment always wins with
$\sigma(10+u)\approx 1$, so the network becomes context-degenerate rather
than better-routed. A four-arm control settles what the bias pays for:
holding the gate at the base system's own measured constant reproduces the
base system, while deleting the gate reproduces most of the repair's gain and
is numerically above it on Fashion-MNIST, a difference that is not significant at
$n{=}10$. The bias removes a constant
damping factor that a never-used gate had left in place.

So the honest summary is this: we measured the size and the cost of the
repair precisely, we bounded where it applies, and we located its mechanism
only far enough to know it is not the one we first assumed. The gap we
closed is real and useful, but it is not the gap we thought we were closing,
and closing it by restoring routing would have required a routing mechanism
this architecture does not have.

Two practical consequences follow. For anyone measuring segment routing in
this family: the dominant-segment instrument must be paired with a mortality
probe, and the two zero-conditions (a never-positive post-kWTA activity rate
and an identically-zero gate row) are not the same statistic --- our own
cross-architecture comparison is narrower than it might appear for exactly
that reason. For anyone proposing a repair of this shape: check whether it
works by fixing the mechanism or by deleting it, and in the latter case ask
what was being deleted. Both are cheap to check --- one control arm, one
probe --- and both would have caught our own misreading earlier than they
did.

Two results from the added controls change the paper's
position relative to the method it audits, and we state them without
discretion. First, the granularity result does not survive a
construction that is not self-limiting: a segment-resolved anchor with
refreshed reference states gains about two and a half points of learning term over
the neuron-resolved one on MNIST, where our original construction
produced no difference at all --- at a cost of roughly twenty points of
final accuracy (Sec.~\ref{sec:ctrl-anchor}). The general claim that
sub-neuron protection granularity is worthless is not supported by our
data and is withdrawn; what the control shows is that the direction of
the effect is set by the reference-update rule, not by the granularity
itself. Second, and more bluntly, DeLP is not competitive with
contemporary protection methods: a replay buffer with logit distillation
reaches $86.4$ against DeLP's $64.0$ on Permuted-MNIST and $75.3$
against $42.5$ on Permuted-Fashion, with markedly less forgetting in
both (Sec.~\ref{sec:ctrl-baselines}). The right way to read this paper is
therefore not as a defence of DeLP but as an argument that the
dendritic-segment account of continual learning needs to be rethought
from the trunk upward --- the routing it assumes does not emerge, the
instrument used to check it is unreliable, the repair we proposed turns
out to be a damping-factor removal, and the whole scheme is outcompeted
by a buffer.

Finally, the meta-result is worth stating plainly. The repair that this
manuscript originally reported as causal survived a $22$-cell
dose--response, a $160$-run transfer test, and a $320$-run factor scan
before three controls took it apart. Each of those controls was cheap. We
had the instrumentation to run them, and we did not run them first. Auditing
the assumption before building on it took one paper's worth of null results;
we hope the corrected instrumentation and the failure modes documented here
make the next audit shorter.

% =====================================================================

\section*{Data availability}
All code, per-run records (JSON), analysis scripts, pre-registration documents
and their dated amendments, and the complete set of superseded records (with
an explicit deprecation notice) are available in the accompanying artifact
repository (URL de-anonymized for double-blind review; the archival DOI
of each record is given in the bibliography). This manuscript merges
three previously separate studies; the corresponding experiment series are
\texttt{s1\_*}, \texttt{gevo\_*}, \texttt{c2*}, \texttt{S3\_CNN},
\texttt{S3b\_E4}, and \texttt{S4\_gate\_bypass}. Figures are produced by the
archived scripts (\texttt{paper\_s1/make\_figs.py} and the S3/S4 equivalents).

\section*{Declaration of competing interest}
The author declares no competing interests.

\bibliographystyle{tmlr}
\bibliography{refs}

\end{document}
